\documentclass[12pt, oneside]{book}

%% Language and font encodings
\usepackage[english]{babel}
\usepackage[utf8x]{inputenc}
\usepackage[T1]{fontenc}
\usepackage{amsfonts}

%% Sets page size and margins
\usepackage[a4paper,top=3cm,bottom=2cm,left=3cm,right=3cm,marginparwidth=1.75cm]{geometry}

%% Useful packages
\usepackage{amsmath}
\usepackage{amssymb}
\usepackage{stmaryrd}
\usepackage{amsthm}
\usepackage{mathrsfs}
\usepackage{graphicx}
\usepackage[colorinlistoftodos]{todonotes}
\usepackage[colorlinks=true, allcolors=black]{hyperref}
\usepackage{enumitem}
\setlist[enumerate,1]{label={(\alph*)}}

%Commands definitions
\newcommand{\Iint}[2]{\llbracket #1 , #2 \rrbracket}
\newcommand{\Zn}[1]{\mathbb{Z} / #1 \mathbb{Z}}
\newcommand{\td}{\textcolor{red}{\textbf{TODO}}}
\newcommand{\A}{\mathbb{A}}
\newcommand{\Z}{\mathbb{Z}}
\newcommand{\Q}{\mathbb{Q}}
\newcommand{\R}{\mathbb{R}}
\newcommand{\C}{\mathbb{C}}
\renewcommand{\O}{\mathcal{O}}
\newcommand{\m}{\mathfrak{m}}
\renewcommand{\P}{\mathbb{P}}

%Command execution. 
%If this line is commented, then the appearance remains as usual.
%\invertbackgroundtext

% New mathematical operator
\DeclareMathOperator{\Ima}{im}
\DeclareMathOperator{\Ker}{Ker}
\DeclareMathOperator{\lcm}{lcm}
\DeclareMathOperator{\sgn}{sgn}
\DeclareMathOperator{\rad}{Rad}

%Set QED symbol to blacksquare
\renewcommand\qedsymbol{$\blacksquare$}

%% Table packages
\usepackage{array}
\newcolumntype{P}[1]{>{\centering\arraybackslash}p{#1}}
\usepackage{multirow}

% Définir un nouveau compteur d'exercice
\newcounter{exercise}[chapter] % Compteur qui est remis à zéro à chaque nouvelle section

% Définir un environnement pour l'exercice
\newenvironment{exercise}{
    \refstepcounter{exercise} % Incrémenter le compteur
    \noindent\textbf{\arabic{chapter}.\arabic{exercise}.}
}{} % Fin de l'environnement

% Définir un environnement pour les solutions
\newenvironment{solution}{
    \noindent\textbf{Solution} 
}{}


\title{Solutions to Algebraic Curves - An Introduction to Algebraic Geometry - William Fulton}
\author{Samy Lahlou}

\begin{document}
\maketitle

\chapter*{Preface}

The goal of this document is to share my personal solutions to the exercises
in the book Algebraic Curves - An Introduction to Algebraic Geometry by William Fulton during my reading. \\
As a disclaimer, the solutions are not unique and there will probably be better or more optimized solutions than mine. Feel free to correct me or ask me anything about the content of this document at the 
following address :

\noindent samy.lahloukamal@mcgill.ca

\tableofcontents

%CHAPTER 1
\include{CHAP 1/chap1}

%CHAPTER 2
\chapter{Affine Varieties}

\section{Coordinate Rings}

\begin{exercise}
    Show that the map that associates to each $F \in k[X_1, ..., X_n]$ a polynomial function in $\mathcal{F}(V,k)$ is a ring homomorphism whose kernel is $I(V)$. \\
\end{exercise}

\begin{solution}
    Let $\varphi : k[X_1, \dots, X_n] \to \mathcal{F}(V,k)$ be the map that sends each polynomial to its associated function on $V$. Let's show that it is a ring homomorphism. It is clear that $F(\lambda) = \lambda$ for all $\lambda \in k$. Moreover, since $\varphi$ is simply a restriction map, then it preserves addition and multiplication of polynomials. Thus, it is a ring homomorphism.
    
    Next, suppose that $F \in \ker \varphi$, i.e $\varphi(F) = 0$, then it follows that $F(x) = 0$ for all $x \in V$, and hence, $F \in I(V)$. Conversely if $F \in I(V)$, then $F(x) = 0$ for all $x \in V$ so $\varphi(F) = 0$ and hence, $F \in \ker \varphi$. Therefore, $\ker \varphi = I(V)$. \\
\end{solution}

\begin{exercise}
    Let $V \subset \A^n$ be a variety. A \textit{subvariety} of $V$ is a variety $W \subset \A^n$ that is contained in $V$. Show that there is a natural one-to-one correspondence between algebraic subsets (resp. subvarieties, resp. points) of $V$ and radical ideals (resp. prime ideals, resp. maximal ideals) of $\Gamma(V)$. (See Problems $1.22$, $1.38$.)\\
\end{exercise}

\begin{solution}
    This problem is strictly the same as Problem 1.38. \\
\end{solution}

\begin{exercise}
    Let $W$ be a subvariety of a variety $V$, and let $I_V(W)$ be the ideal of $\Gamma(V)$ corresponding to $W$. (a) Show that every polynomial function on $V$ restricts to a polynomial function on $W$. (b) Show that the map from $\Gamma(V)$ to $\Gamma(W)$ defined in part (a) is a surjective homomorphism with kernel $I_W(V)$, so that $\Gamma(W)$ is isomorphic to $\Gamma(V)/I_V(W)$.\\
\end{exercise}

\begin{solution}
    \begin{enumerate}[label=(\alph*)]
        \item Let $f \in \mathcal{F}(V,k)$ be a polynomial function, then there is a polynomial $F \in k[X_1, ..., X_n]$ such that $F(x) = f(x)$ for all $x \in V$. Hence, the function $f|_W \in \mathcal{F}(W,k)$ is also a polynomial function since $f|_W(x) = f(x) = F(x)$ for all $x \in W$.
        \item Let $\varphi : \Gamma(V) \to \Gamma(W)$ be the map that sends a polynomial function on $V$ to its restriction on $W$. This is a ring homomorphism because, clearly $\varphi(\lambda) = \lambda$ for all $\lambda \in k$, and because it is a restriction map so it preserves addition and multiplication.
        
        It is a surjective homomorphism because given any function $f \in \Gamma(W)$, we know that there exists a polynomial $F \in k[X_1, ..., X_n]$ such that $f(x) = F(x)$ for all $x \in W$. Hence, if we define $g \in \Gamma(V)$ by $g(x) = F(x)$ for all $x \in V$, then clearly $\varphi(g) = f$. It follows that $\varphi$ is surjective.
        
        If $f \in \ker(\varphi)$, then it is equivalent to say that $f$ is sent to 0 by $\varphi$. Equivalently, it means that $f(x) = 0$ for all $x \in W$. But by definition, this is equivalent to $f \in I_V(W)$ since $I_V(W)$ is the projection of $I(V)$ in $\Gamma(V)$. Thus, $\ker(\varphi) = I_V(W)$. Therefore, by the First Isomorphism Theorem: $\Gamma(V)/ I_V(W) = \Gamma(W)$.  \\
    \end{enumerate}
\end{solution}

\begin{exercise}
    Let $V \subset \A^n$ be a nonempty variety. Show that the following are equivalent: (i) $V$ is a point; (ii) $\Gamma(V) = k$; (iii) $\dim_k\Gamma(V) < \infty$. \\
\end{exercise}

\begin{solution}
    If $V$ is a point $(a_1, ..., a_n)$, then $I(V) = (X_1 - a_1, ..., X_n - a_n)$ which implies that $\Gamma(V) = k[X_1, ..., X_n] / (X_1 - a_1, ..., X_n - a_n) = k$. Hence, (i) implies (ii).
    
    Next, if $\Gamma(V) = k$, then clearly $\dim_k \Gamma(V) = \dim_k k = 1 < \infty$. Hence, (ii) implies (iii).
    
    Finally, if $\dim_k \Gamma(V) < \infty$, then $V(I(V)) = V$ is a finite set. But since $V$ is irreducible, then it must be a point. Thus, (iii) implies (i). Therefore, the three propositions are equivalent. \\
\end{solution}

\begin{exercise}
    Let $F$ be an irreducible polynomial in $k[X,Y]$, and suppose $F$ is monic in $Y$: $F = Y^n + a_1(X)Y^{n-1} +\dots + a_n(X)$, with $n > 0$. Let $V = V(F) \subset \A^2$. Show that the natural homomorphism from $k[X]$ to $\Gamma(V) = k[X,Y]/(F)$ is one-to-one, so that $k[X]$ may be regarded as a subring of $\Gamma(V)$; show that the residues $\overline{1}, \overline{Y}, ..., \overline{Y}^{n-1}$ generate $\Gamma(V)$ over $k[X]$ as a module. \\
\end{exercise}

\begin{solution}
    Both propositions follow from the fact that $\Gamma(V) = k[X,Y]/(F) = k[X][Y]/(F) \cong k[X][Y]_{< n}$ which denotes the ring of polynomials in $Y$ of degree at most $n-1$ (addition and multiplication modulo $F$) with coefficients in $k[X]$. Hence, $k[X]$ is simply the subring of "constant" polynomials; and the elements $\overline{1}, \overline{Y}, ..., \overline{Y}^{n-1}$ generate $\Gamma(V)$ because the elements in $\Gamma(V)$ are polynomials of $Y$ of degree at most $n-1$.
\end{solution}
\section{Polynomial Maps}

\begin{exercise}
    Let $\varphi : V \to W$, $\psi : W \to Z$. Show that $\widetilde{\psi \circ \varphi} = \tilde{\phi} \circ \tilde{\psi}$. Show that the composition of polynomial maps is a polynomial map. \\
\end{exercise}

\begin{solution}
    For all functions $f \in \mathcal{F}(Z,k)$, we have
    $$(\widetilde{\psi \circ \varphi})(f) = f \circ \psi \circ \varphi = \tilde{\varphi}(f \circ \psi) = \tilde{\varphi}(\tilde{\psi}(f)) = (\tilde{\phi} \circ \tilde{\psi})(f).$$
    Thus, $\widetilde{\psi \circ \varphi} = \tilde{\phi} \circ \tilde{\psi}$.
    
    Next, suppose that $V \subset \A^n$, $W \subset \A^m$ and $Z \subset \A^k$. If $\varphi$ and $\psi$ are polynomial maps, then there exist polynomials $T_1, ..., T_m \in k[X_1, ..., X_m]$ and polynomials $F_1, ...$, $F_k$ $ \in k[X_1, ..., X_m]$ such that $\varphi(x) = (T_1(x), ..., T_m(x))$ for all $x \in \A^n$ and $\psi(x) = (F_1(x), ..., F_k(x))$ for all $x \in \A^m$. It follows that
    $$(\psi\circ \varphi)(x) = \psi(T_1(x), ..., T_m(x)) = (F_1(T_1(x), ..., T_m(x)), ..., F_k(T_1(x), ..., T_m(x)))$$
    for all $x \in \A^n$. For all $i = 1, ..., k$, the map that sends $x \in \A^n$ to $F_i(T_1(x), ..., T_m(x))$ is a polynomial because the composition of polynomials is also a polynomial. There-fore, the composition of two polynomial maps is a polynomial map.\\ 
\end{solution}

\begin{exercise}
    If $\varphi : V \to W$ is a polynomial map, and $X$ is an algebraic subset of $W$, show that $\varphi^{-1}(X)$ is an algebraic subset of $V$. If $\varphi^{-1}(X)$ is irreducible, and $X$ is contained in the image of $\varphi$, show that $X$ is irreducible. This gives a useful test for irreducibility. \\
\end{exercise}

\begin{solution}
    Let $X = V(F_1, ..., F_r)$, then $x \in \varphi^{-1}(X)$ iff $\varphi(x) \in X$, iff $F_i(\varphi(x)) = 0$ for all $i = 1, ..., r$. Therefore, $\varphi^{-1}(X) = V(F_1\circ \varphi, ..., F_r \circ \varphi)$ which shows that $\varphi^{-1}(X)$ is an algebraic set ($F_i \circ \varphi$ is a polynomial for all $i$). Suppose now that $\varphi^{-1}(X)$ is irreducible, and $X$ is contained in the image of $\varphi$. If $X = X_1 \cup X_2$ where $X_1, X_2$ are algebraic sets such that $X_i$ is not contained in $X_j$ for $i \neq j$, then $\varphi^{-1}(X) = \varphi^{-1}(X_1) \cup \varphi^{-1}(X_2)$. By the first part of this exercise, we have that both $\varphi^{-1}(X_1)$ and $\varphi^{-1}(X_2)$ are algebraic sets. Moreover, the fact that $X$ is contained in the image of $\varphi$ implies that $\varphi^{-1}(X_i) \subset \varphi^{-1}(X_j)$ for $i \neq j$ can never happen. Thus, $\varphi^{-1}(X)$ is not irreducible, a contradiction. Therefore, by contradiction, $X$ is irreducible. \\
\end{solution}

\begin{exercise}
    (a) Show that $\{(t,t^2, t^3) \in \A^3(k) \ | \ t\in k\}$ is an affine variety. (b) Show that $V(XZ - Y^2, YZ - X^3, Z^2 - X^2Y) \subset \A^3(\C)$ is a variety. (\textit{Hint:} $Y^3 - X^4, Z^3 - X^5, Z^4 - Y^5 \in I(V)$. Find a polynomial map from $\A^1(\C)$ onto $V$.) \\
\end{exercise}

\begin{solution}
    \begin{enumerate}[label=(\alph*)]
        \item Problem 1.33.
        \item Define the polynomial map $\varphi: \A^1(\C) \to \A^3(\C)$ that maps a complex number $t$ to $(t^3, t^4, t^5)$. It is easy to verify that $\varphi(t) \in V(I)$ for all $t \in \C$ (which implies that $\varphi^{-1}(V(I)) = \A^1(\C)$). Let's show that $V(I)$ is in the image of $\varphi$. If $(x,y,z) \in V(I)$, then $xz = y^2$, $yz = x^3$, and $z^2 = x^2y$. If $x = 0$ or $y = 0$, then it is clear that $(x,y,z) = (0,0,0) \in \Ima \varphi$. Otherwise, we have that $z = y^2/x$ and $z = x^3/y$ which implies that $y^2/x = x^3/y$, and hence, $y^3 = x^4$. If we let $t$ be a cube root of $x$, then $y^3 = t^{12}$. Since $t$ can be any cube root of $x$, then we can choose, in particular the cube root such that $y = t^4$. Finally, since $yz = x^3$, then $z = t^9/t^4 = t^5$. Thus, $(x,y,z) = \varphi(t)$ which implies that $V(I) \subset \Ima \varphi$. 
        
        Next, since $\varphi^{-1}(V(I)) = \A^1(\C)$ and $\A^1(\C)$ is irreducible, then by Problem 2.7, $V(I)$ is irreducible. \\
    \end{enumerate}
\end{solution}

\begin{exercise}
    Let $\varphi: V \to W$ be a polynomial map of affine varieties, $V' \subset V$, $W' \subset W$ subvarieties. Suppose $\varphi(V') \subset W'$. (a) Show that $\tilde{\varphi}(I_W(W')) \subset I_V(V')$ (see Problem 2.3). (b) Show that the restriction of $\varphi$ gives a polynomial map from $V'$ to $W'$. \\
\end{exercise}

\begin{solution}
    \begin{enumerate}[label=(\alph*)]
        \item Let $\tilde{\varphi}(f) = f \circ \varphi \in \tilde{\varphi}(I_W(W'))$ with $f \in I_W(W')$, then for all $x \in V'$, $\varphi(x) \in W'$ which implies that $(f\circ \varphi)(x) = 0$ since $f \in I_W(W')$. It follows that $\tilde{\varphi}(f) \in I_V(V')$, and hence, $\tilde{\varphi}(I_W(W')) \subset I_V(V')$.
        \item Follows directly from the fact that $\varphi(V') \subset W'$, and the fact that $\varphi$ is a polynomial map. \\
    \end{enumerate}
\end{solution}

\begin{exercise}
    Show that the \textit{projection map} $\text{pr} = \A^n \to \A^r$, $n \geq r$, defined by $\text{pr}(a_1, ..., a_n) = (a_1, ..., a_r)$ is a polynomial map. \\ 
\end{exercise}

\begin{solution}
    If we let $P_i \in k[X_1, ..., X_n]$ be the polynomial $P_i(X_1, ..., X_n) = X_i$ for all $i = 1, ..., r$, then clearly $\text{pr} = (P_1, ..., P_r)$. \\
\end{solution}

\begin{exercise}
    Let $f \in \Gamma(V)$, $V$ a variety $\subset \A^n$. Define
    $$G(f) = \{(a_1, ..., a_n, a_{n+1}) \in \A^{n+1} \ | \ (a_1, ..., a_n) \in V \text{ and } a_{n+1} = f(a_1, ..., a_n)\},$$
    the \textit{graph} of $f$. Show that $G(f)$ is an affine variety, and that the map $(a_1, ..., a_n) \mapsto (a_1, ..., a_n, f(a_1, ..., a_n))$ defines an isomorphism of $V$ with $G(f)$. (Projection gives the inverse.) \\
\end{exercise}

\begin{solution}
    If $V = V(F_1, ..., F_r)$, then $G(f) = V(F_1, ..., F_r, X_{n+1} - f(X_1, ..., X_n))$; hence, it is an algebraic set. Next, consider the map $\varphi: V \to \A^{n+1}$ that maps $(a_1, ..., a_n)$ to $(a_1, ..., a_n, f(a_1, ..., a_n))$. This is a polynomial map with image equal to $G(f)$. Since $\varphi^{-1}(G(f)) = V$ is irreducible, then it follows that $G(f)$ is irreducible by Problem 2.7.
    
    Next, let's show that $\varphi$ is an isomorphism. Let $\psi : G(f) \to V$ be the projection map. This is a polynomial map by Problem 2.10. Let $v \in V$, then $\psi(\varphi(v)) = \psi(v, f(v)) = v$; let $(x, f(x)) \in G(f)$, then $\varphi(\psi(x, f(x))) = \varphi(x) = (x, f(x))$. Therefore, $\varphi$ is an isomorphism. \\
\end{solution}

\begin{exercise}
    (a) Let $\varphi : \A^1 \to V = V(Y^2 - X^3) \subset \A^2$ be defined by $\varphi(t) = (t^2, t^3)$. Show that although $\varphi$ is a one-to-one, onto polynomial map, $\varphi$ is not an isomorphism. (\textit{Hint:} $\tilde{\varphi}(\Gamma(V)) = k[T^2, T^3] \subset k[T] = \Gamma(\A^1)$.) (b) Let $\varphi : \A^1 \to V = V(Y^2 - X^2(X+1))$ be defined by $\varphi(t) = (t^2-1, t(t^2-1))$. Show that $\varphi$ is one-to-one and onto, except that $\varphi(\pm 1) = (0,0)$.\\
\end{exercise}

\begin{solution}
    \begin{enumerate}[label=(\alph*)]
        \item Clearly, $\varphi$ is a polynomial map. If $\varphi(s) = \varphi(t)$, then $s^2 = t^2$ and $s^3 = t^3$. If $t = 0$, then $s = 0$ and so $t = s$. Otherwise, since $s^2 = t^2$, then $s = \pm t$. If $s = -t$, then cubing both sides gives us $s^3 = -t^3$ which implies that $t^3 = -t^3$, and hence, $t = 0$, a contradiction. It follows that $s = -t$ cannot hold, and hence, $s = t$. Therefore, $\varphi$ is one-to-one. To show that $\varphi$ is onto, let $(x,y) \in V$, then $y^2 = x^3$. If we let $t$ be a square root of $x$, then we get that $y^2 = t^6$. It follows that $y = \pm t^3 = (\pm t)^3$. If $y = t^3$, then we are done since $(x,y) = \varphi(t)$; if $y = (-t)^3$, then we get that $(x,y) = \varphi(-t)$. It follows that $\varphi$ is onto.
    
        Suppose that there is a polynomial map $\psi : V \to \A^1$ such that $\psi \circ \varphi = id$, then $\tilde{\varphi} \circ \tilde{\psi} = \tilde{id} = id : k[X] \to k[X]$ by Problem 2.6 using the fact that $ \Gamma(\A^1) = k[X]$. But this implies that $\Ima \tilde{\varphi} = k[X]$ which is impossible since $\Ima \tilde{\varphi} = k[X^2, X^3] \not \ni X$. Therefore, $\psi$ cannot exist, and $\varphi$ cannot be an isomorphism.
        \item If $\varphi(s) = \varphi(t)$, then $s^2 - 1 = t^2 - 1$ and $s(s^2 - 1) = t(t^2 - 1)$. Suppose that $s^2 - 1 = t^2 - 1 \neq 0$, then $s(s^2 - 1) = t(t^2 - 1)$ directly implies that $s = t$. The case $s^2 - 1 = t^2 - 1 = 0$ implies that $s = \pm 1$ and $t = \pm 1$ which we don't consider here. Next, let $(x,y) \in V$, then $y^2 = x^2(x+1)$. If we let $t$ be a root of the equation $t^2 - 1 = x$, then $y^2 = (t^2 - 1)^2(t^2 - 1 +1) = [t(t^2 - 1)]^2$. It follows that $y = \pm t(t^2 - 1)$. Equivalently, $(x,y) = \varphi(\pm t)$ which, in both cases, implies that $(x,y) \Ima \varphi$. Therefore, $\varphi$ is one-to-one and onto, except at $\pm 1$.
    \end{enumerate}
\end{solution}

\begin{exercise}
    Let $V = V(X^2 - Y^3, Y^2 - Z^3) \subset \A^3$ as in Problem 1.40, $\overline{\alpha} : \Gamma(V) \to k[T]$ induced by the homomorphism $\alpha$ of that problem. (a) What is the polynomial map $f$ from $\A^1$ to $V$ such that $\tilde{f} = \overline{\alpha}$? (b) Show that $f$ is one-to-one and onto, but not an isomorphism. \\
\end{exercise}

\begin{solution}
    First, recall that $\alpha : k[X,Y,Z] \to k[T]$ is defined by $\alpha(X) = T^9$, $\alpha(Y) = T^6$, and $\alpha(Z) = T^4$. Since $\ker \alpha = I(V)$, then we have an induced homomorphism $\overline{\alpha} : \Gamma(V) \to k[T]$ such that $\overline{\alpha} \circ \pi = \alpha$ where $\pi : k[X,Y,Z] \to \Gamma(V)$ is the projection map. 
    \begin{enumerate}[label=(\alph*)]
        \item By following the proof of Proposition 1 (or simply by looking at the definition of $\alpha$), we get that $f : \A^1 \to V$ defined by $f(T) = (T^9, T^6, T^4)$ is the polynomial map that satisfies $\tilde{f} = \varphi$.
        \item If $f(s) = f(t)$, then $s^k = t^k$ for $k = 4,6,9$. We can assume that $s,t \neq 0$ since otherwise, $s = t$. From $s^4 = t^4$, we get that $s = i^mt$ for some integer $m = 0,1,2,3$ where $i$ denotes a root of $-1$. Taking both sides to the sixth power gives us $s^6 = (-1)^mt^6$. But since $s^6 = t^6$, then $t^6 = (-1)^mt^6$ which implies that $(-1)^m = 1$, and hence, that $m = 0,2$. It follows that $s = (-1)^pt$ for some integer $p = 0,1$. Next, taking the previous equation to the power of nine gives us that $s^9 = (-1)^pt^9$; but since $s^9 = t^9$, then $t^9 = (-1)^pt^9$, and hence $(-1)^p = 1$. It follows that $s = t$. Thus, $f$ is one-to-one.
        
        Let $(x,y,z) \in V$, then $x^2 = y^3$ and $y^2 = z^3$. If we let $t$ be a fourth root of $z$, then $y^2 = t^{12}$ which implies that $y = \pm t^6$, and hence, that $y = (i^kt)^6$ where $i$ is a root of $\sqrt{-1}$, and $k$ is any of $0,2$ or $1,3$. Similarly, if we do the same with the ninth power, we get that $x = (i^kt)^9$ for some $k$ which will correspond to the case of the sixth power. It follows that $(x,y,z) = f(i^k t)$, and hence, $f$ is onto $V$. Finally, let's show that $f$ is not an isomorphism. Suppose that there is a polynomial map $g : V \to \A^1$ such that $g \circ f = id$, then $\tilde{f} \circ \tilde{g} = \tilde{id} = id : k[T] \to k[T]$ by Problem 2.6 using the fact that $ \Gamma(\A^1) = k[T]$. But this implies that $\Ima \tilde{f} = k[T]$ which is impossible since $\Ima \tilde{f} = k[T^9, T^6, T^4] \not \ni T$. Therefore, $g$ cannot exist, and hence, $f$ cannot be an isomorphism. 
    \end{enumerate}
\end{solution}
\section{Coordinate Changes}

\begin{exercise}
    A set $V \subset \A^n(k)$ is called a \textit{linear subvariety} of $\A^n(k)$ if $V = V(F_1, ..., F_r)$ for some polynomials $F_i$ of degree 1. (a) Show that if $T$ is an affine change of coordinates on $\A^n$, then $V^T$ is also a linear subvariety of $\A^n(k)$. (b) If $V \neq \varnothing$, show that there is an affine change of coordinates $T$ of $\A^n$ such that $V^T = V(X_{m+1}, ..., X_n)$. (\textit{Hint:} use induction on $r$.) So $V$ is a variety. (c) Show that the $m$ that appears in part (b) is independent of the choice of $T$. It is called the dimension of $V$. Then $V$ is then isomorphic (as a variety) to $\A^m(k)$. (\textit{Hint:} Suppose there were an affine change of coordinates $T$ such that $V(X_{m+1}, ..., X_n)^T = V(X_{s+1}, ..., X_n)$, $m < s$; show that $T_{m+1}, ..., T_n$ would be dependent.) \\
\end{exercise}

\begin{solution}
    \begin{enumerate}[label=(\alph*)]
        \item If we denote $I = (F_1, ..., F_r)$, then $I^T$ is generated by the elements $F^T$ with $F \in I$. But notice that if $F = \sum_i a_iF_i$, then
        $$F^T = \tilde{T}(\sum_i a_iF_i) = \sum_i \tilde{T}(a_i)\tilde{T}(F_i) = \sum_i a_i^T F_i^T$$
        which implies that $I^T = (F_1^T, ..., F_r^T)$. Since the $F_i$'s and the $T_i$'s are all linear, then the $F_i^T$ are linear, and hence, $V^T = V(F_1^T, ..., F_r^T)$ is a linear subvariety.
        \item Let's prove it by induction on $r\geq 1$. When $r = 1$, we have that $V = V(f)$ where $\deg(f) = 1$. If we write $f = a_0 + \sum_i a_iX_i$, then we must have at least one $i$ such that $a_i \neq 0$, call it $i = i_0$. Define the polynomial map $T = (T_1, ..., T_n)$ where $T_i = 0$ for all $i \neq i_0$, and $T_{i_0} = a_{i_0}^{-1}(X_n - a_0)$, then
        $$V^T = V(f^T) = V\left(a_0 + \sum_i a_iT_i\right) = V\left(a_0 + a_{i_0}[a_{i_0}^{-1}(X_n - a_0)]\right) = V(X_n).$$
        Hence, the base case holds. Next, suppose that the statement holds for some $r = k$ and consider the case where $V = V(f_1, \dots, f_k, f_{k+1})$. By the induction hypothesis, there is a polynomial $T$ such that
        $$V^T = V(f_1^T, \dots, f_k^T)\cap V(f_{k+1}^T) = V(X_{m+1}, ..., X_n, f_{k+1}^T).$$
        Now, notice that if we write $f = a_0 + \sum_{i=1}^{n}a_i X_i$, then $(X_{m+1}, ..., X_n, f_{k+1}^T) = (X_{m+1}, ..., X_n, g)$ where $g = a_0 + \sum_{i=1}^{m}a_i X_i$. Hence: $V^T = V(g, X_{m+1}, ..., X_n)$. If $g = 0$, we are done, if $g \neq 0$, then there is a $i = i_0 \leq m$ such that $a_{i_0} \neq 0$. Thus, if we define the polynomial map $U = (U_1, ..., U_n)$ by $U_i = 0$ for all $1\leq i\neq i_0 \leq m$, $U_{i_0} = a_{i_0}^{-1}(X_m - a_0)$, and $U_i = X_i$ for all $m+1 \leq i \leq n$. Thus, by construction:
        $$V^{T\circ U} = V(g\circ U, X_{m+1}\circ U, ..., X_n \circ U) = V(X_m, X_{m+1}, ..., X_n).$$
        Thus, the case $r = k+1$ also holds. Therefore, by induction, the proposition is proved for all $r \geq 1$.
        \item Suppose there exist affine change of coordinates $U_1$ and $U_2$ such that
        $$V^{U_1} = V(X_{m+1}, \dots, X_n) \qquad \text{ and } \qquad V^{U_2} = V(X_{s+1}, \dots, X_n)$$
        with $m < s$, then equivalently, there is a change of coordinates $T$ such that $V(X_{m+1}, ..., X_n)^T = V(X_{s+1}, ..., X_n)$. If we write $T = (T_1, ..., T_n)$, we get that $V(T_{m+1}, ..., T_n) = V(X_{s+1}, ..., X_n)$. By the Nullstellensatz, and the fact that $(X_{s+1}, ..., X_n)$ is a prime ideal:
        $$\{T_{m+1}, ..., T_n\} \subseteq \rad(T_{m+1}, ..., T_n) = I(V(T_{m+1}, ..., T_n)) = (X_{s+1}, ..., X_n).$$
        Thus, for all $m+1 \leq i \leq n$, $T_i \in (X_{s+1}, ..., X_n)$. Equivalently, 
        $$T_i = \alpha_{s+1, i}X_{s+1} + \dots + \alpha_{n,i}X_n$$
        for some scalars $\alpha_{j,i} \in k$, and for all $m+1 \leq i \leq n$ (technically, the coefficients should be polynomials but we can deduce that there must be such scalars by the fact that $T_i$ has degree 1). Now, recall that $T$ is an affine change of coordinates so its linear part must be invertible. Equivalently, the matrix form of its linear part must be invertible, and hence, its rows must be linearly independent vectors. But the above equation shows that the last $n -m$ rows of that matrix are of the form $(0,..., 0, \alpha_{s+1,i}, ..., \alpha_{n,i})$. In other words, the vectors
        $$\begin{pmatrix}
            0 \\ \vdots \\ 0 \\ \alpha_{s+1,m+1} \\ \vdots \\ \alpha_{n,m+1}
        \end{pmatrix}, \qquad \begin{pmatrix}
            0 \\ \vdots \\ 0 \\ \alpha_{s+1,m+2} \\ \vdots \\ \alpha_{n,m+2}
        \end{pmatrix}, \qquad \dots, \qquad \begin{pmatrix}
            0 \\ \vdots \\ 0 \\ \alpha_{s+1,n} \\ \vdots \\ \alpha_{n,n}
        \end{pmatrix}$$
        are linearly independent. This is equivalent to the fact that the vectors
        $$\begin{pmatrix}
            \alpha_{s+1,m+1} \\ \vdots \\ \alpha_{n,m+1}
        \end{pmatrix}, \qquad \begin{pmatrix}
            \alpha_{s+1,m+2} \\ \vdots \\ \alpha_{n,m+2}
        \end{pmatrix}, \qquad \dots, \qquad \begin{pmatrix}
            \alpha_{s+1,n} \\ \vdots \\ \alpha_{n,n}
        \end{pmatrix}$$
        are linearly. But this is impossible since we have $n - m$ vectors in $k^{n - s}$, and $n - m > n-s$. Therefore, we found a contradiction that proves that $m$ is independent of the choice of $T$.
        
        We now need to prove that $V$ is isomorphic to $\A^m(k)$ as a variety. Let $T$ be an affine coordinates change such that $V^T = V_m := V(X_{m+1}, ..., X_n)$, then $T(V_m) = V$. It follows that $T$, restricted to $V_m$, is a bijection with image equal to $V$. Since the inverse of $T$ is also an affine change of coordinates, then it follows that $V$ and $V_m$ are isomorphic as variety. Thus, we now need to show that $V_m$ is isomorphic to $\A^m(k)$ as variety. Consider the polynomial map $U: \A^m(k) \to V_m$ defined by
        $$U(X_1, ..., X_m) = (X_1, ..., X_m, 0, ..., 0).$$
        Its inverse is the polynomial map $\pi : V_m \to \A^m(k)$ defined by
        $$\pi(X_1, ..., X_m, X_{m+1}, ..., X_n) = (X_1, ..., X_m).$$
        By construction, they are inverses of each other so $U$ is an isomorphism of variety. Thus, $V$ is isomorphic to $V_m$, which is in turn isomorphic to $\A^m(k)$ as a variety. \\
    \end{enumerate}
\end{solution}

\begin{exercise}
    Let $P = (a_1, ..., a_n)$, $Q = (b_1, ..., b_n)$ be distinct points in $\A^n$. The \textit{line} through $P$ and $Q$ is defined to be $\{a_1 + t(b_1 - a_1), ..., a_n + t(b_n - a_n) | t\in k\}$. (a) Show that if $L$ is the line through $P$ and $Q$, and $T$ is an affine change of coordinates, then $T(L)$ is the line through $T(P)$ and $T(Q)$. (b) Show that a line is a linear subvariety of dimension 1, and that a linear subvariety of dimension 1 is the line through any of its two points. (c) Show that, in $\A^2$, a line is the same thing as a hyperplane. (d) Let $P,P' \in \A^2$, $L_1,L_2$ two distinct lines through $P$, $L_1', L_2'$ are distinct lines through $P'$. Show that there is an affine change of coordinates $T$ of $\A^2$ such that $T(P) = P'$ and $T(L_i) = L_i'$, $i=1,2$. \\
\end{exercise}

\begin{solution}
    \begin{enumerate}
        \item First, rewrite the line through $P$ and $Q$ as $L = \{P + t(Q-P) | t \in k\}$. If we let $T = T' + \alpha$ be an affine change of coordinates where $T'$ is the linear part, and $\alpha \in k^n$ is a constant, then 
        \begin{align*}
            T(L) &= \{T(P + t(Q-P)) | t \in k\} \\
            &= \{T'(P + t(Q-P)) + \alpha | t \in k\} \\
            &= \{T'(P) + t(T'(Q) - T'(P)) + \alpha | t \in k\} \\
            &= \{[T'(P) + \alpha] + t([T'(Q) + \alpha] - [T'(P) + \alpha]) | t \in k\} \\
            &= \{T(P) + t(T(Q) - T(P)) | t \in k\}
        \end{align*}
        which is precisely the definition of the line through $T(P)$ and $T(Q)$.
        \item First, consider the affine change of coordinates $T_1$ that maps $x$ to $x - P$ (a translation), then the image of $L$ is the line $L_1$ through 0 with director vector $v_1 = Q-P$. Next, since $v_1$ is a non-zero vector, then we can find vectors $v_2, ..., v_n$ such that $v_1, ..., v_n$ forms a basis of $k^n$. From that, define the linear transformation $T_2$ that maps $v_i$ to $e_i$ where $e_1, ..., e_n$ is the standard basis. By construction, $T_2$ is an invertible linear transformation, so it is an affine change of coordinates. The image of $L_1$ under $T_2$ is the line $L_2 = V(X_2, ..., X_n)$, which is a linear subvariety. Thus, by construction, we have that $L = T_1^{-1}(T_2^{-1}(V)) = V^{T_2 \circ T_1}$ is also a linear subvariety (by Problem 2.14 part (a)).
        
        Conversely, let $L$ be a linear subvariety of dimension 1, then there is an affine change of coordinates $T$, such that $L^T = V(X_2, ..., X_n)$, or equivalently, such that $L = T(V(X_2, ..., X_n))$. Let $P$ and $Q$ be two points in $L$, then there are two points $(p, 0, ..., 0), (q, 0, ..., 0) \in V(X_2, ..., X_n)$ that are the respective preimages of these two points. Since $P$ and $Q$ are distinct, and $T$ is a bijection, then $p$ and $q$ are also distinct. Hence, we can easily see that $V(X_2, ..., X_n)$ can be described as the line $L_0$ through the vectors $(p, 0, ..., 0)$ and  $(q, 0, ..., 0)$. Finally, using part (a), we get that $L = T(L_0)$ must be the line through the points $P$ and $Q$. Since $P \neq Q$ were arbitrary points on $L$, then we can say that a linear subvariety fo dimension 1 is the line through any two of its points.
        \item Let $L$ be a line in $\A^2$, then by the previous part, we know that there exists an affine change of coordinates $T$ such that $L^T = V(X_2)$, or equivalently, such that $L = T(V(X_2))$. Since $T$ is a bijection such that its inverse is also an affine change of coordinates, we can rewrite the previous equation as
        $$L = T(V(X_2)) = (T^{-1})^{-1}(V(X^2)) = V(X_2)^{T^{-1}}.$$
        Hence, if we write $T^{-1} = (T_1, T_2)$, we get that $L = V(X_2 \circ T^{-1}) = V(T_2)$. Therefore, $L$ is a hyperplane ($T_2$ has degree 1). 

        Conversely, let $V(f)$ be a hyperplane, then $f$ is a polynomial of degree 1. By the proof of the base case of the induction in the solution of part (b) in Problem 2.14, we have that there exists an affine change of coordinates $T$ such that $V(f)^T = V(X_2)$. Since the latter is the line through $(0,0)$ and $(1,0)$, and $T(V(X_2)) = V(f)$, then $V(f)$ must be a line. Therefore, in $\A^2$, hyperplanes and lines form the same class of objects.
        \item First, define the affine change of coordinates $T_1$ that acts as a translation by mapping the point $P$ to the origin. Thus, the image of the lines $L_1$ and $L_2$ are now intersecting at the origin. Let $v_1$ and $v_2$ be director vectors of the lines $T(L_1)$ and $T(L_2)$, since the lines are distinct, the vectors form a basis of $\A^2$. Similarly, repeat the same process for the lines $L_1'$ and $L_2'$: let $T_2$ be an affine change of coordinates that maps $P'$ to the origin (a translation), and let $v_3$, $v_4$ be the director vectors of the lines $T_2(L_1')$ and $T_2(L_2')$ respectively. Again, $v_3$, $V_4$ forms another basis of $\A^2$. Finally, consider the linear transformation that maps $v_1$ to $v_3$ and $v_2$ to $v_4$, then it is invertible, and so it is an affine change of coordinates. By construction, $T(T_1(L_1)) = T_2(L_1')$ and $T(T_1(L_2)) = T_2(L_2')$. Therefore, the affine change of coordinates $U = T_2^{-1}\circ T \circ T_1$ has the property that $U(P) = P'$ and $U(L_i) = L_i'$ for $i = 1,2$. \\
    \end{enumerate}
\end{solution}

\begin{exercise}
    Let $k = \C$. Give $\A^n(\C) = \C^n$ the usual topology (obtained by identifying $\C$ with $\R^2$, and hence $\C^n$ with $\R^{2n}$). Recall that a topological space $X$ is \textit{path-connected} if for any $P,Q \in X$, there is a continuous mapping $\gamma : [0,1] \to X$ such that $\gamma(0) = P$ and $\gamma(1) = Q$. (a) Show that $\C \setminus S$ is path-connected for any finite set $S$. (b) Let $V$ be an algebraic set in $\A^n(\C)$. Show that $\A^n(\C) \setminus V$ is path-connected. (\textit{Hint:} If $P,Q \in \A^n(\C) \setminus V$, let $L$ be the line through $P$ and $Q$. Then $L \cap V$ is finite, and $L$ is isomorphic to $\A^1(\C)$.) \\
\end{exercise}

\begin{solution}
    \begin{enumerate}
        \item Let $S = \{s_1, ..., s_n\}$ and consider the line segment $L$ between $P$ and $Q$. If $L\cap S = \varnothing$, then it suffices to take $\gamma(t) = p + t(q-p)$. Otherwise, take any point $l \in L$ and consider the straightline $L_0$ through $l$ that is perpendicular to $L$. For all points $z \in L_0$, we can define the line segments $L_1^{(z)}$ between $P$ and $z$, $L_2^{(z)}$ between $z$ and $Q$. For two distinct $z_1, z_2 \in L_0$, we have that
        $$(L_1^{(z_1)} \cup L_2^{(z_1)})\cap (L_1^{(z_2)} \cup L_2^{(z_2)}) = \{P,Q\}.$$
        Thus, if we define $L_z = L_1^{(z)} \cup L_2^{(z)}$, then there must be a $z_0 \in L_0$ such that $L_z \cap S = \varnothing$ since there are infinitely many points on $L_0$ but only finitely many in $S$. From that, define 
        $$\gamma(t) = \begin{cases}
            p+2t(z_0 - p) & 0 \leq t < \frac{1}{2}, \\
            z_0+2t(q - z_0) & \frac{1}{2} \leq t \leq 1. \\
        \end{cases}.$$
        After easily verifying that $\gamma$ is continuous, we are done.
        \item Let $P,Q \in \A^n(\C) \setminus V$ and $L = \{P + t(Q-P) | t \in \C\}$ be the line through $P$ and $Q$. If we write $V = V(f_1, ..., f_r)$, then
        $$L\cap V = (L\cap V(f_1)) \cap V(f_2, ..., f_r) \subseteq L\cap V(f_1).$$
        If we look at $L \cap V(f_1)$ more closely, we get that it is the set of solutions of the polynomial $f_1$ such that the input is in the line $L$. If we define $g(t) = P + t(Q-P)$, we get that the elements are precisely the image under $g$ of the solutions of the function $f_1(g(t))$. But notice now that $f_1 \circ g$ is only one variable so it has finitely many solutions. Thus, $L \cap V(f_1)$ is finite, and hence, so is $L \cap V$.
        
        Next, since $L$ is a linear subvariety of dimension 1, then it is isomorphic to $\A^1(\C)$ as a variety (Problem 2.14 part (c)). Hence, there is a bijective polynomial map $T : \C \to L$. Let $S = T^{-1}(L \cap V)$, then $S$ is finite. By part (a), there is a path $\gamma : [0,1] \to \C\setminus S$ from $P' = T^{-1}(P)$ to $Q' = T^{-1}(Q)$. It follows that $T \circ \gamma : [0,1] \to L \setminus V \subseteq \A^n(\C)\setminus V$ is path from $P$ to $Q$. Therefore, $\A^n(\C)$ is path-connected. \\
    \end{enumerate}
\end{solution}
\section{Rational Functions and Local Rings}

\begin{exercise}
    Let $V = V(Y^2 - X^2(X+1)) \subset \A^2$, and $\overline{X}$, $\overline{Y}$ the residues of $X$, $Y$ in $\Gamma(V)$; let $z = \overline{Y}/\overline{X} \in k(V)$. Find the pole sets of $z$ and of $z^2$. \\
\end{exercise}

\begin{solution}
    First, notice that $z^2 = \overline{X+1}$ so its pole set is the empty set. Let's prove that $(0,0)$ is the only element in the pole set of $z$. Since $z = \overline{Y}/\overline{X}$ and $z = \overline{X(X+1)}/\overline{Y}$, then $(0,0)$ is the only possible element in the pole set. If the pole set is empty, then $z$ must be a rational function, and hence, we have the equality $\overline{Y}/\overline{X} = \overline{p(X,Y)}$ which is equivalent to
    $$Y - Xp(X,Y) = q(X,Y)(Y^2 - X^2(X+1))$$
    in $k[X,Y]$. Taking $X=0$ gives us the equation $Y = q(0,Y)Y^2$ which is impossible. Therefore, the pole set of $z$ is composed of $(0,0)$ only.\\
\end{solution}

\begin{exercise}
    Let $\mathcal{O}_P(V)$ be the local ring of a variety $V$ at a point $P$. Show that there is a natural one-to-one correspondence between the prime ideals in $\mathcal{O}_P(V)$ and the subvarieties of $V$ that pass through $P$. (\textit{Hint:} If $I$ is prime in $\mathcal{O}_P(V)$, $I\cap \Gamma(V)$ is prime in $\Gamma(V)$, and $I$ is generated by $I\cap \Gamma(V)$; use Problem 2.2.)\\
\end{exercise}

\begin{solution}
    By Problem 2.2, we know that there is a natural one-to-one correspondence between the subvarieties of $V$ and the prime ideals of $\Gamma(V)$. If $I$ is a prime ideal of $\mathcal{O}_P(V)$, then $I \cap \Gamma(V)$ is a prime ideal of $\Gamma(V)$. Conversely, if $I$ is a prime ideal of $\Gamma(V)$, then viewed as a subset of $\mathcal{O}_P(V)$, it generates a prime ideal $I'$ of $\mathcal{O}_P(V)$ (same idea as in the proof of Proposition 3). This establishes a natural one-to-one correspondence between the prime ideals of $\Gamma(V)$ and the prime ideals of $\mathcal{O}_P(V)$. Therefore,  there is a natural one-to-one correspondence between the prime ideals in $\mathcal{O}_P(V)$ and the subvarieties of $V$. \\
\end{solution}

\begin{exercise}
    Let $f$ be a rational function on a variety $V$. Let $U = \{P \in V\ | \ f$ is defined at $p\}$. Then $f$ defines a function from $U$ to $k$. Show that this function determines $f$ uniquely. So a rational function may be considered as a type of function, but only on the complement of an algebraic subset of $V$, not on $V$ itself. \\
\end{exercise}

\begin{solution}
    Let $f$ and $g$ be two rational functions in $k(V)$ that have the same induced function, then $h = f-g$ is a rational function in $k(V)$ such that its induced function is the zero function on its domain. Hence, it suffices to show that if $h \in k(V)$ induces the zero function on its domain $U$, then it is the zero rational function. To do so, write $h = \overline{a}/\overline{b}$ where $\overline{a}, \overline{b} \in \Gamma(V)$ and such that $b(P) \neq 0$ for some $P \in V$ (i.e, $h$ is defined on at least one point). By our assumption on $h$, we have that $a(P)/b(P) = 0$ whenever $h$ is defined at $p$, or equivalently, whenever $b(P) \neq 0$. But notice that $a(P)/b(P) = 0$ implies that $a(P) = 0$. Hence, $b(P) \neq 0$ implies that $a(P) = 0$. It follows that the polynomial $ab$ satisfies $a(P)b(P) = 0$ for all $P \in V$. Thus, $ab \in I(V)$, and hence, $\overline{a}\overline{b} = 0$ in $\Gamma(V)$. But since $V$ is a variety, $\Gamma(V)$ is an integral domain so we must have $\overline{a} = 0$ or $\overline{b} = 0$. Since we know that $\overline{b}$ is non-zero on at least one point of $V$, then we get that $\overline{a} = 0$. It follows that $h = 0$ as a rational function. Therefore, $f = g$. \\
\end{solution}

\begin{exercise}
    In the example given in this section, show that it is impossible to write $f = a/b$, where $a,b \in \Gamma(V)$, and $b(P) \neq 0$ for every $P$ where $f$ is defined. Show that the pole set of $f$ is exactly $\{(x,y,z,w) | y=0 \text{ and } w=0\}$. \\
\end{exercise}

\begin{solution}
    First, suppose by contradiction that we can write $f = \overline{a}/\overline{b}$ where
    $$b(x_0,y_0,z_0,w_0) = 0 \implies y_0 = w_0 = 0,$$
    then $b(x_0,y_0,z_0,w_0) \neq 0$ whenever $y_0 \neq 0$ or $w_0 \neq 0$. Hence, if $y_0 \neq 0$, the polynomial $b(x,y_0, z, w)$ has no roots, and hence, it is equal to a non-zero constant. Similarly, when $w_0 \neq 0$, the polynomial $b$ is constant. Let $(x_i,y_i,z_i,w_i) \in \A^4\setminus V(y,w)$ for $i=1,2$. Without loss of generality, we can assume that $y_1, w_2 \neq 0$. Thus, we get that
    $$b(x_1, y_1, z_1, w_1) = b(x_2, y_1, z_2, w_2) = b(x_2, y_2, z_2, w_2) =: c \in \A^1\setminus \{0\}$$
    using our previous remarks. This proves that $b$ is constant on the set $\A^4\setminus V(y,w)$. Now, fix $(x_0,z_0, w_0) \in \A^4$ and consider the polynomial $b_0(y) = b(x_0, y, z_0, w_0)$. By what we proved, we know that this polynomial is constant on the set $\A^1\setminus\{0\}$. Our goal is to show that $b_0(0) = c$, and hence, that $b_0$ is constant over its domain. If $b_0(0) \neq 0$, then $b_0$ has no roots and so it is a constant polynomial. It follows that $b_0 \equiv c$. If $b_0(0) = 0$, then $b_0(y) = y^m b_1(y)$ where $m$ is the multiplicity of the root, and $b_1$ is a polynomial that doesn't vanish at $y = 0$. Since the roots of $b_1$ are necessarily roots of $b_0$, then $b_1$ has no roots. Thus, it is a non-zero constant polynomial, giving us $b_0(y) = Cy^m$ where $C\neq 0$. But notice that $k$ is algebraically closed so there is certainly an $m$th root of 2 which we will call $\alpha$. Thus, $b_0(1) = C$ and $b_0(\alpha) = 2C$, contradicting the fact that $b_0$ is constant on $\A^1\setminus \{0\}$. Hence, the only possible case is that $b_0 \equiv c$. Since it holds for all $(x_0, z_0, w_0) \in \A^3$, we get that $b \equiv c$. Thus, $f = \frac{1}{c}\overline{a} \in \Gamma(V)$. Let's now prove that this is impossible. Notice that $f = \overline{x}/\overline{y}$ implies that
    $$x = \frac{1}{c}ya(x,y,z,w) + h(x,y,z,w)(xw - yz).$$
    Taking $y = w = 0$ gives us $x = 0$, a contradiction. Therefore, by contradiction, such a polynomial $b$ cannot exist.\\

    Next, let's prove that the pole set of $f$ is exactly $V(y,w)$. From the two expressions of $f$: $f = \overline{x}/\overline{y} = \overline{z}/\overline{w}$, we get that the pole set must be a subset of $V(y,w)$. Conversely, let $p = (p_x, 0, p_z, 0) \in V(y,w)$ and suppose by contradiction that we can write $f = \overline{a}/\overline{b}$ where $b(p) \neq 0$. From the equation $\overline{x}/\overline{y} = \overline{a}/\overline{b}$, we get
    $$xb(x,y,z,w) = ya(x,y,z,w) + h_1(x,y,z,w)(xw - yz).$$
    Taking $y = w = 0$ gives us
    $$xb(x,0,z,0) = 0.$$
    But this is impossible because we know that $b(x,0,z,0)$ is a non-zero polynomial ($b(p_x, 0, p_z, 0) \neq 0$). Thus, by contradiction, $p$ must be in the pole set of $f$. It follows that the pole set of $f$ is $V(y,w)$. \\
\end{solution}

\begin{exercise}
    Let $\varphi : V \to W$ be a polynomial map of affine varieties, $\tilde{\varphi} : \Gamma(W) \to \Gamma(V)$ the induced map on coordinate rings. Suppose $P \in V$, $\varphi(P) = Q$. Show that $\tilde{\varphi}$ extends uniquely to a ring homomorphism (also written $\tilde{\varphi}$) from $\mathcal{O}_Q(W)$ to $\mathcal{O}_P(V)$. (Note that $\tilde{\varphi}$ may not extend to all of $k(W)$.) Show that $\tilde{\varphi}(\mathfrak{m}_Q(W)) \subset \mathfrak{m}_P(V)$. \\
\end{exercise}

\begin{solution}
    Define $\psi : \O_Q(W) \to \O_P(V)$ by $\psi(a/b) = \tilde{\varphi}(a)/\tilde{\varphi}(b) = (a \circ \varphi)/(b \circ \varphi)$. First, let's show that its maps to $\O_P(V)$. Let $a/b \in O_Q(W)$ such that $b(Q) \neq 0$, then $(b\circ \varphi)(P) = b(Q) \neq 0$. It follows that $\tilde{\varphi}(a) / \tilde{\varphi}(b) \in O_P(V)$. Next, we need to show that the map is well-defined. Suppose that $a/b = c/d \in O_Q(W)$, then $ad = bc$ as elements of $\Gamma(W)$. Since $\tilde{\varphi}$ is a morphism on $\Gamma(W)$, we get that $\tilde{\varphi}(a)\tilde{\varphi}(d) = \tilde{\varphi}(b)\tilde{\varphi}(c)$, and hence, $\tilde{\varphi}(a)/\tilde{\varphi}(b) = \tilde{\varphi}(c)/\tilde{\varphi}(d)$. Therefore, $\psi$ is a well-defined map from $\O_Q(W)$ to $\O_P(V)$.\\
    
    Now, we need to show that $\psi$ is a ring homomorphism:
    $$\psi(1) = \psi(1/1) = \tilde{\varphi}(1)/\tilde{\varphi}(1) = 1/1 = 1,$$
    and
    \begin{align*}
        \psi((a/b) + (c/d)) &= \psi((ad + bc)/bd) \\
        &= \tilde{\varphi}(ad+bc)/\tilde{\varphi}(bd) \\
        &= (\tilde{\varphi}(a)\tilde{\varphi}(d) + \tilde{\varphi}(b)\tilde{\varphi}(c))/\tilde{\varphi}(b)\tilde{\varphi}(d) \\
        &= (\tilde{\varphi}(a)/\tilde{\varphi}(b)) + (\tilde{\varphi}(c)/\tilde{\varphi}(d)) \\
        &= \psi(a/b) + \psi(c/d),
    \end{align*}
    and 
    \begin{align*}
        \psi((a/b)\cdot (c/d)) &= \psi((ac)/(bd)) \\
        &= \tilde{\varphi}(ac)/\tilde{\varphi}(bd) \\
        &= (\tilde{\varphi}(a)\tilde{\varphi}(c))/(\tilde{\varphi}(b)\tilde{\varphi}(d)) \\
        &= (\tilde{\varphi}(a) / \tilde{\varphi}(b))\cdot (\tilde{\varphi}(c)/\tilde{\varphi}(d)) \\
        &= \psi(a/b)\psi(c/d).\\
    \end{align*}
    Next, we need to show that it extends $\tilde{\varphi}$. For all $a \in \Gamma(W)\subseteq O_Q(W)$, we have
    $$\psi(a) = \psi(a/1) = \tilde{\varphi}(a)/\tilde{\varphi}(1) = \tilde{\varphi}(a)/1 = \tilde{\varphi}(a).$$

    Now, let $\psi_0 : \O_Q(W)\to \O_P(V)$ be a ring homomorphism that extends $\tilde{\varphi}$, then for all $a/b$ in $\O_Q(W)$:
    \begin{align*}
        & \tilde{\varphi}(b)\psi_0(a/b) = \psi_0(b)\psi_0(a/b) = \psi_0(b\cdot (a/b)) = \psi_0(a) = \tilde{\varphi}(a)\\
        &\implies \psi_0(a/b) = \tilde{\varphi}(a)/\tilde{\varphi}(b) = \psi(a/b).
    \end{align*}
    Therefore, $\psi_0 = \psi$, and hence, $\psi$ is unique. Thus, we will now define $\tilde{\varphi}$ to be $\psi$.\\

    Finally, let $a/b \in \m_Q(W)$, then $a(Q) = 0$ which gives us
    $$\tilde{\varphi}(a/b)(P) = (a\circ \varphi)(P)/(b\circ \varphi)(P) = a(Q)/b(Q) = 0$$
    so $\tilde{\varphi}(a/b) \in \m_P(V)$. Thus, $\tilde{\varphi}(\m_Q(W)) \subset \m_P(V)$. \\
\end{solution}

\begin{exercise}
    Let $T : \A^n \to \A^n$ be an affine change of coordinates, $T(P) = Q$. Show that $\tilde{T} : \O_Q(\A^n) \to \O_P(\A^n)$ is an isomorphism. Show that $\tilde{T}$ induces an isomorphism from $\O_Q(V)$ to $\O_P(V^T)$ if $P \in V^T$, for $V$ a subvariety of $\A^n$. \\
\end{exercise}

\begin{solution}
    Since $T$ is an isomoprhism from $\A^n$ to $\A^n$, then $\tilde{T} : \Gamma(\A^n) \to \Gamma(\A^n)$ is also an isomorphism with inverse $\tilde{T}^{-1} : \Gamma(\A^n) \to \Gamma(\A^n)$. By Problem 2.21, both $\tilde{T}$ and $\tilde{T}$ must act by
    $$\tilde{T}(a/b) = \tilde{T}(a)/\tilde{T}(b), \quad \text{ and } \quad \tilde{T}^{-1}(c/d) = \tilde{T}^{-1}(c)/\tilde{T}^{-1}(d).$$
    Hence, for all $a/b \in \O_Q(\A^n)$ and $c/d \in \O_P(\A^n)$:
    $$\tilde{T}^{-1}(\tilde{T}(a/b)) = \tilde{T}^{-1}(\tilde{T}(a))/\tilde{T}^{-1}(\tilde{T}(b)) = a/b$$
    and 
    $$\tilde{T}(\tilde{T}^{-1}(c/d)) = \tilde{T}(\tilde{T}^{-1}(c))/\tilde{T}(\tilde{T}^{-1}(d)) = c/d.$$
    Therefore, $\tilde{T} : \O_Q(\A^n) \to \O_P(\A^n)$ is an isomorphism. The proof that $\tilde{T}$ induces an isomorphism from $\O_Q(V)$ to $\O_P(V^T)$ is the same as the one we just did since $T$ is an isomorphism from $V^T$ to $V$. \\
\end{solution}
\section{Discrete Valuation Rings}

\begin{exercise}
    Show that the order function on $K$ is independent of the choice of uniformizing parameter. \\
\end{exercise}

\begin{solution}
    Let $t_1$ and $t_2$ be two uniformizing parameters for $R$. Since $t_2$ is clearly non-zero, then we must have that $t_2 = ut_1^{n}$ for some unit $u$ and $n \geq 1$. Similarly, we also have $t_1 = vt_2^{m}$ for some unit $v$ and $m \geq 1$. It follows that $t_2 = uv^nt_2^{nm}$. By the uniqueness of the expressions of the form $u_0t_2^{n_0}$, we get that $nm = 1$, and hence, that $n = 1$. It follows that $t_2 = ut_1$ for some unit $u$.
    
    Next, let $z \in R\setminus\{0\}$, then we have that $z = u_1t_1^{n_1}$ and $z = u_2 t_2^{n_2}$ for some units $u_1,u_2$ and integers $n_1, n_2 \geq 0$. Since $t_2 = ut_1$, then $u_1t_1^{n_1} = u_2t_2^{n_2}$ implies that
    $$u_1t_1^{n_1} = (u_2u^{n_2})t_1^{n_2}.$$
    Since $u_2u^{n_2}$ is unit, we get that $n_1 = n_2$ by uniqueness. Therefore, the order of $z$ doesn't depend on the uniformizing parameter. \\
\end{solution}

\begin{exercise}
    Let $V = \A^1$, $\Gamma(V) = k[X]$, $K = k(V) = k(X)$. (a) For each $a \in k = V$, show that $\O_a(V)$ is a DVR whith uniformizing parameter $t = X-a$. Show that $\O_{\infty} = \{F/G \in k(X) \ | \ \deg(G) \geq \deg(F)\}$ is also a DVR, with uniformizing parameter $t = 1/X$. \\
\end{exercise}

\begin{solution}
    Since $k[X]$ is a UFD, every element in $\O_a(V) \subset k(X)$ has a unique representation $P/Q$ where $P$ and $Q$ have no common factors. By properties of polynomials of one variable, we can write $P(X) = (X-a)^m P_0(X)$ where $m \geq 0$ and such that $P_0(a) \neq 0$. Hence, $P/Q = (P_0/Q)(X-a)^m$. The uniqueness of this expression comes from the unique factorization in $k[X]$. The rational function $P_0/Q$ is a unit since $Q/P_0 \in \O_a(V)$. Since the element $t = X-a$ is irreducible, then $\O_a(V)$ is a DVR with uniformizing parameter $t = X-a$.
    
    By the unique factorization in $k[X]$, we can write each element $F/G$ in $\O_{\infty}$ as $(F_0/G_0)\cdot (1/X)^m$ where $F_0$ and $G_0$ are non divisible by $X$, and where $m \in \Z$. Hence, $\deg(G_0) + m \geq \deg(F_0)$. Next, define $k = \deg(F_0) - \deg(G_0)$, then $(m-k) + \deg(G_0) + k \geq \deg(F_0)$ and $\deg(G_0) + k = \deg(F_0)$. From this, we get that we can write 
    $$\frac{F}{G} = \left(\frac{F_0}{G_0}\cdot X^{-k}\right)\cdot \left(\frac{1}{X}\right)^{m-k}$$
    If $m - k < 0$, then
    $$\deg(F_0) \leq \deg(G_0) + m < \deg(G_0) + k = \deg(F_0)$$
    which is a contradiction. Thus, $m - k \geq 0$. Moreover, the term $(F_0/G_0)X^{-k}$ is a unit in $\O_{\infty}$ because the degree of the numerator is equal to the degree of the denominator. Next, suppose that we can write
    $$\frac{1}{X} = \frac{P_1}{Q_1}\cdot \frac{P_2}{Q_2}$$
    where $(P_1/Q_1), (P_2/Q_2) \in \O_{\infty}$ and none are units, then $\deg(P_1) < \deg(Q_1)$ and $\deg(P_2) < \deg(Q_2)$. Equivalently, $\deg(P_1) + 1 \leq \deg(Q_1)$ and $\deg(P_2) + 1 \leq \deg(Q_2)$. Hence:
    $$2 + \deg(P_1) + \deg(P_2) \leq \deg(Q_1) + \deg(Q_2).$$
    Now, if we rearrange the above equation, we obtain
    $$Q_1 Q_2 = XP_1P_2$$
    from which we get
    \begin{align*}
        \deg(Q_1) + \deg(Q_2) &= 1 + \deg(P_1) + \deg(P_2) \\
        &< 2 + \deg(P_1) + \deg(P_2) \\
        &\leq \deg(Q_1) + \deg(Q_2),
    \end{align*}
    a contradiction. Hence, $1/X$ is irreducible in $\O_{\infty}$. Finally, suppose that
    $$\frac{F_0}{G_0}\left(\frac{1}{X}\right)^{n_0} = \frac{F_1}{G_1}\left(\frac{1}{X}\right)^{n_1}$$
    where $(F_0/G_0), (F_1/G_1) \in \O_{\infty}$ are units, then $\deg(F_0) = \deg(G_0)$ and $\deg(F_1) = \deg(G_1)$. If we rearrange the equation, we get
    $$F_0G_1X^{n_1} = F_1G_0X^{n_0}$$
    which gives us
    $$\deg(F_0) + \deg(G_1) + n_1 = \deg(F_1) + \deg(G_0) + n_0.$$
    Hence, $n_0 = n_1$, and $F_0/G_0 = F_1/G_1$. Therefore, the expressions of that form are unique, and so $\O_{\infty}$ is a DVR with uniformizing paramter $t = 1/X$. \\
\end{solution}

\begin{exercise}
    Let $p \in \Z$ be a prime number. Show that $\{r \in \Q \ | \ r = a/b, \ a,b \in \Z, \ p$ doesn't divide $b\}$ is a DVR with quotient field $\Q$. \\
\end{exercise}

\begin{solution}
    Let's first show that set we are interested in, which we will call $R$, is a DVR. First, $R$ is a ring because it is a subset of the ring $\Q$ that contains the additive and multiplicative identities, and that is closed under addition and multiplication. Moreover, it is a domain since it is a subring of a domain. Let $a/b \in R$, then we can write $a = a_0 p^n$ where $n \geq 0$ and $p$ doesn't divide $a_0$. Hence, $a/b = (a_0/b)p^n$ and $a_0/b$ is a unit since $b/a_0 \in R$. 
    
    Next, let $a_i/b_i \in R^{\times}$ and $n_i \geq 0$ for $i = 0,1$, then none of $a_0, a_1, b_0, b_1$ are divisible by $p$. Suppose that
    $$\frac{a_0}{b_0}p^{n_0} = \frac{a_1}{b_1}p^{n_1},$$
    then 
    $$a_0b_1p^{n_0} = a_1b_0 p^{n_1}.$$
    By the Fundamental Theorem of Arithmetic, $n_0 = n_1$, and hence, $a_0/b_0 = a_1/b_1$. Therefore, the expressions of that form are unique. Finally, the element $p$ is irre-ducible in $R$ because if $p = (a_0/b_0)(a_1/b_1)$ where both factors are non units, then $a_i = pc_i$ for $i = 0,1$. Hence, $pb_0b_1 = p^2 c_0c_1$ which is impossible since the left hand side is only divisible by $p$ once. Thus, $p$ is irreducible, and so $R$ is a DVR with uniformizing parameter $t = p$.
    
    The last thing we need to show is that the quotient field of $R$, denoted by $\text{Frac}(R)$, is $\Q$. If we recall that the quotient field of a ring is the smallest field that contains it, then $\Z \subseteq R$ implies that $\Q = \text{Frac}(\Z) \subseteq \text{Frac}(R)$ and $R \subseteq \Q$ implies that $\text{Frac}(R) \subseteq \Q$. Thus, $\text{Frac}(R) = \Q$. \\
\end{solution}

\begin{exercise}
    Let $R$ be a DVR with quotient field $K$; let $\m$ be the maximal ideal of $R$. (a) Show that if $z \in K$, $z \notin R$, then $z^{-1} \in \m$. (b) Suppose $R \subset S \subset K$, and $S$ is also a DVR. Suppose the maximal ideal of $S$ contains $\m$. Show that $S = R$. \\
\end{exercise}

\begin{solution}
    \begin{enumerate}
        \item Since $R$ is a DVR, then every non-zero element of $K$ can be written as $ut^n$ where $u \in R^{\times}$, $t$ is a fixed uniformizing parameter for $R$, and $n \in \Z$. Hence, if $z \in K \setminus R$, then we must have that $z \neq 0$ and hence, that $z = ut^n$ with $n < 0$. It follows that $z^{-1} = u^{-1}t^{-n}$. Since $u^{-1} \in R^{\times}$ and $-n > 0$, then $z^{-1} = R \setminus R^{\times} = \m$.
        \item First, let $t$ be a uniformizing parameter for $S$ and suppose that $t \notin R$, then $t \in S\setminus R \subset K \setminus R$. Thus, by part (a), $t^{-1} \in \m \subset R \subset S$. But this is impossible since $t$ is not a unit in $S$. Thus, $t \in R$. 
        
        Next, let's show that $t$ is a uniformizing parameter for $R$. Since $t \in R$ and $R$ is a DVR, then $t = ut_0^{n_0}$ where $t_0$ is a uniformizing parameter for $R$, $u \in R^{\times} \subset S^{\times}$ and $n_0 \geq 0$. Similarly, $t_0 \in R\setminus\{0\} \subset S\setminus\{0\}$ so $t_0 = vt^{n_1}$ where $v \in S^{\times}$ and $n_1 \geq 0$. It follows that $t = (uv^{n_0})t^{n_0n_1}$. Viewing this equation as an equation in $S$, we get that $n_0n_1 = 1$, and hence, $n_0 = 1$. Thus, $t = ut_0$ where $u \in R^{\times}$. Therefore, $t$ is a uniformizing parameter for $R$.

        Finally, let $z$ be an arbitrary non-zero element in $S$, then $z = ut^n$ where $u$ is a unit, and $n \geq 0$. Since $t \in R$, then $t^n \in R$. Suppose that $u \notin R$, then $u \in K \setminus R$, and hence, $u^{-1} \in \m$. Since $t$ is a uniformizing parameter in $R$ and $u^{-1}$ is not a unit in $R$, then $u^{-1} = vt^m$ for some unit $v \in R^{\times}$ and $m \geq 1$. Bt if we view this equation as an equation in $S$, we get that $m = 0$, a contradiction. So we must have $u \in R$ as well, and hence, $z = ut^n$. Therefore, $S = R$. 
        
        [\textbf{This proof never uses the assumption that the maximal ideal of $S$ contains $\m$. After rereading my proof, it seems like there is no mistake, and hence, the assumption seems to be unecessary.}]\\
    \end{enumerate}
\end{solution}

\begin{exercise}
    Show that the DVR's of Problem 2.24 are the only DVR's with quotient field $k(X)$ that contain $k$. Show that those of Problem 2.25 are the only DVR's with quotient field $\Q$. \\
\end{exercise}

\begin{solution}
    \td  \\
\end{solution}

\begin{exercise}
    An \textit{order function} on a field $K$ is a function $\varphi$ from $K$ onto $\Z \cup \{\infty\}$, satisfying:
    \begin{enumerate}
        \item[(i)] $\varphi(a) = \infty$ if and only if $a = 0$.
        \item[(ii)] $\varphi(ab) = \varphi(a) + \varphi(b)$.
        \item[(iii)] $\varphi(a + b) \geq \min(\varphi(a), \varphi(b))$.  
    \end{enumerate}
    Show that $R = \{z \in K | \varphi(z) \geq 0\}$ is a DVR with maximal ideal $\m = \{z | \varphi(z) > 0\}$, and quotient field $K$. Conversely, show that if $R$ is a DVR with quotient field $K$, then the function $\text{ord}: K \to \Z\cup \{\infty\}$ is an order function on $K$. Giving a DVR with quotient field $K$ is equivalent to defining an order function on $K$. \\
\end{exercise}

\begin{solution}
    \td \\
\end{solution}

\begin{exercise}
    \td \\
\end{exercise}

\begin{solution}
    \td \\
\end{solution}

\begin{exercise}
    \td \\
\end{exercise}

\begin{solution}
    \td \\
\end{solution}

\begin{exercise}
    \td \\
\end{exercise}

\begin{solution}
    \td \\
\end{solution}

\begin{exercise}
    \td \\
\end{exercise}

\begin{solution}
    \td \\
\end{solution}


\section{Forms}

\begin{exercise}
    \td \\
\end{exercise}

\begin{solution}
    \td \\
\end{solution}

\begin{exercise}
    \td \\
\end{exercise}

\begin{solution}
    \td \\
\end{solution}

\begin{exercise}
    \td \\
\end{exercise}

\begin{solution}
    \td \\
\end{solution}

\begin{exercise}
    \td \\
\end{exercise}

\begin{solution}
    \td \\
\end{solution}
\include{CHAP 2/SECTION 2.7/section2.7}
\section{Operations with Ideals}

\begin{exercise}
    \td \\
\end{exercise}

\begin{solution}
    \td \\
\end{solution}

\begin{exercise}
    \td \\
\end{exercise}

\begin{solution}
    \td \\
\end{solution}

\begin{exercise}
    \td \\
\end{exercise}

\begin{solution}
    \td \\
\end{solution}

\begin{exercise}
    \td \\
\end{exercise}

\begin{solution}
    \td \\
\end{solution}

\begin{exercise}
    \td \\
\end{exercise}

\begin{solution}
    \td \\
\end{solution}

\begin{exercise}
    \td \\
\end{exercise}

\begin{solution}
    \td \\
\end{solution}

\begin{exercise}
    \td \\
\end{exercise}

\begin{solution}
    \td \\
\end{solution}

\begin{exercise}
    \td \\
\end{exercise}

\begin{solution}
    \td \\
\end{solution}
\section{Ideals with a Finite Number of Zeros}

\begin{exercise}
    \td \\
\end{exercise}

\begin{solution}
    \td \\
\end{solution}
\section{Quotient Modules and Exact Sequences}

\begin{exercise}
    \td \\
\end{exercise}

\begin{solution}
    \td \\
\end{solution}

\begin{exercise}
    \td \\
\end{exercise}

\begin{solution}
    \td \\
\end{solution}

\begin{exercise}
    \td \\
\end{exercise}

\begin{solution}
    \td \\
\end{solution}

\begin{exercise}
    \td \\
\end{exercise}

\begin{solution}
    \td \\
\end{solution}

\begin{exercise}
    \td \\
\end{exercise}

\begin{solution}
    \td \\
\end{solution}

\begin{exercise}
    \td \\
\end{exercise}

\begin{solution}
    \td \\
\end{solution}
\section{Free Modules}

\begin{exercise}
    \td \\
\end{exercise}

\begin{solution}
    \td \\
\end{solution}

\begin{exercise}
    \td \\
\end{exercise}

\begin{solution}
    \td \\
\end{solution}

\begin{exercise}
    \td \\
\end{exercise}

\begin{solution}
    \td \\
\end{solution}

%CHAPTER 3
\chapter{Local Properties of Plane Curves}

\section{Multiple Points and Tangent Lines}

\begin{exercise}
    Prove that in the above examples $P = (0,0)$ is the only multiple point on the curves $C$, $D$, $E$, and $F$. \\
\end{exercise}

\begin{solution}
    \begin{enumerate}
        \item Let $F = Y^2 - X^3$, then $F_X = -3X^2$ and $F_Y = 2Y$. A point $P = (a,b)$ is a multiple point if and only if $F_X(P) = F_Y(P) = 0$. This is equivalent to $-3a^2 = 2b = 0$, i.e., $P = (0,0)$. Therefore, the only multiple point is $P = (0,0)$.
        \item Let $F = Y^2 - X^3 - X^2$, then $F_X = -3X^2 - 2X$ and $F_Y = 2Y$. A point $P = (a,b) \in F$ is a multiple point if and only if $F_X(P) = F_Y(P) = 0$. This is equivalent to $-3a^2 - 2a = 2b = 0$. These equations are equivalent to $a = 0, -2/3$ and $b = 0$, but the point $(-2/3, 0) \notin F$ so the only solution is $P = (0,0)$. Therefore, the only multiple point is $P = (0,0)$. 
        \item Let $F = (X^2 + Y^2)^2 + 3X^2Y - Y^3$, then $F_X = 4X^3 + 4XY^2 + 6XY$ and $F_Y = 4X^2Y + 4Y^3 + 3X^2 - 3Y^2$. A point $P = (a,b) \in F$ is a multiple point if and only if $F_X(P) = F_Y(P) = 0$. This is equivalent to the following system of equations:
        $$\begin{cases}
            4a^3 + 4ab^2 + 6ab = 0 \\
            4a^2b + 4b^3 + 3a^2 - 3b^2 = 0.
        \end{cases}$$
        Clearly, $P = (0,0)$ is a solution of the system and a point in $F$ so it is a multiple point. Moreover, by looking at the definition of $F$, we see that $(0,0)$ is the only point in $F$ such that $X = 0$ or $Y = 0$. Hence, to find the other multiple points, we can assume that there is a multiple point $(a,b)$ such that both $a$ and $b$ are non-zero. Dividing the first equation by $a$ gives us 
        $$4a^2 + 4b^2 + 6b = 0$$
        which is the same as
        $$a^2 = -\frac{1}{2}(2b^2 + 3b).$$
        Replacing $a^2$ in the second equation with the expression above gives us
        $$b = -9/22$$
        after making some simplifications. It follows that $a = 54/121$, and hence, we get the point $(54/121, -9/22)$. However, this point is not in $F$. Therefore, the only multiple point is $P = (0,0)$. 
        \item Let $F = (X^2 + Y^2)^3 - 4X^2Y^2$, then $F_X = 2X[3(X^2 + Y^2)^2 - 4Y^2]$ and $F_Y = 2Y[3(X^2 + Y^2)^2 - 4X^2]$. A point $P = (a,b) \in F$ is a multiple point if and only if $F_X(P) = F_Y(P) = 0$. This is equivalent to the following system of equations:
        $$\begin{cases}
            2a[3(a^2 + b^2)^2 - 4b^2] = 0 \\
            2b[3(a^2 + b^2)^2 - 4a^2] = 0.
        \end{cases}$$
        Clearly, $P = (0,0)$ is a solution of the system and a point in $F$ so it is a multiple point. Moreover, by looking at the definition of $F$, we see that $(0,0)$ is the only point in $F$ such that $X = 0$ or $Y = 0$. Hence, to find the other multiple points, we can assume that there is a multiple point $(a,b)$ such that both $a$ and $b$ are non-zero. Dividing the first equation by $2a$ and the second by $2b$ gives us 
        $$\begin{cases}
            3(a^2 + b^2)^2 - 4b^2 = 0 \\
            3(a^2 + b^2)^2 - 4a^2 = 0.
        \end{cases}$$
        It follows that $4b^2 = 4a^2$, and hence, that $a^2 = b^2$. Replacing $b^2$ with $a^2$ in the first equation of the latter system gives us the equation
        $$3(2a^2)^2 - 4a^2 = 0$$
        which is equivalent to $a = \pm 1/\sqrt{3}$. From $a^2 = b^2$, we get that $b = \pm 1/\sqrt{3}$ as well. However, $F(\pm1/\sqrt{3}, \pm 1\sqrt{3}) \neq 0$ so the points $(\pm 1 / \sqrt{3}, \pm 1/\sqrt{3})$ are not on the curve $F$. Therefore, the only multiple point is $P = (0,0)$. \\
    \end{enumerate}
\end{solution}

\begin{exercise}
    Find the multiple points, and the tangent lines at the multiple points, for each of the following curves:
    \begin{enumerate}
        \item $Y^3 - Y^2 + X^3 - X^2 + 3XY^2 + 3X^2Y + 2XY$
        \item $X^4 + Y^4 - X^2Y^2$
        \item $X^3 + Y^3 - 3X^2 - 3Y^2 + 3XY + 1$
        \item $Y^2 + (X^2 - 5)(4X^4 - 20X^2 + 25)$
    \end{enumerate}
    Sketch the part of the curve in (d) contained in $\A^2(\R) \subset \A^2(\C)$. \\
\end{exercise}

\begin{solution}
    \begin{enumerate}
        \item Let $F = Y^3 - Y^2 + X^3 - X^2 + 3XY^2 + 3X^2Y + 2XY$, then $F_X = 3X^2 - 2X + 3Y^2 + 6XY + 2Y$ and $F_Y = 3Y^2 - 2Y + 6XY + 3X^2 + 2X$. The system $F_X(a,b) = F_Y(a,b) = 0$ is equivalent to \td
        \item \td
        \item \td
        \item \td \\
    \end{enumerate}
\end{solution}

\begin{exercise}
    \td \\
\end{exercise}

\begin{solution}
    \td \\
\end{solution}

\begin{exercise}
    \td \\
\end{exercise}

\begin{solution}
    \td \\
\end{solution}

\begin{exercise}
    \td \\
\end{exercise}

\begin{solution}
    \td \\
\end{solution}

\begin{exercise}
    \td \\
\end{exercise}

\begin{solution}
    \td \\
\end{solution}

\begin{exercise}
    \td \\
\end{exercise}

\begin{solution}
    \td \\
\end{solution}

\begin{exercise}
    \td \\
\end{exercise}

\begin{solution}
    \td \\
\end{solution}

\begin{exercise}
    \td \\
\end{exercise}

\begin{solution}
    \td \\
\end{solution}

\begin{exercise}
    \td \\
\end{exercise}

\begin{solution}
    \td \\
\end{solution}

\begin{exercise}
    \td \\
\end{exercise}

\begin{solution}
    \td \\
\end{solution}
\include{CHAP 3/SECTION 3.2/section3.2}
\include{CHAP 3/SECTION 3.3/section3.3}

%CHAPTER 4
\chapter{Projective Varieties}

\section{Projective Space}

\begin{exercise}
    What points in $\P^2$ do not belong to two of the three sets $U_1, U_2, U_3$? \\
\end{exercise}

\begin{solution}
    By definition, a point $P = [x_1:x_2:x_3] \notin U_i$ if and only if $x_i = 0$. Hence, the point $P$ do not belong to two of the three sets $U_1, U_2, U_3$ if and only if two of its three coordinates are zero. Thus, the possible values of $P$ are $[x : 0 : 0]$, $[0 : y : 0]$, and $[0 : 0 : z]$ where $x,y,z \in k\setminus\{0\}$. Finally, these points are equal to the points $[1 : 0 : 0]$, $[0 : 1 : 0]$, and $[0 : 0 : 1]$. Therefore, these three points are precisely the three points that do not belong to two of the three sets $U_1, U_2, U_3$. \\
\end{solution}

\begin{exercise}
    Let $F \in k[X_1, ..., X_{n+1}]$ ($k$ infinite). Write $F = \sum F_i$, $F_i$ a from of degree $i$. Let $P \in \P^n(k)$, and suppose $F(x_1, ..., x_{n+1}) = 0$ for every choice of homogeneous coordinates $(x_1, ..., x_{n+1})$ for $P$. Show that each $F_i(x_1, ..., x_{n+1}) = 0$ for all homogeneous coordinates for $P$. (\textit{Hint:} consider $G(\lambda) = F(\lambda x_1, \dots, \lambda x_{n+1}) = \sum \lambda^i F_i(x_1, ..., x_{n+1})$ for fixed $(x_1, ..., x_{n+1})$.) \\
\end{exercise}

\begin{solution}
    Let $x \in \A^{n+1}(k)$ be a choice of coordinates for the point $P$, and define the polynomial $G(\lambda) = F(\lambda x) = \lambda^i F_i(x)$. Here, since $x$ is fixed, then $F_i(x)$ can be seen as the coefficients defining $G$. Since $\lambda x$ is another choice of coordinates for the point $P$ for all $\lambda \neq 0$, then $G(\lambda) = F(\lambda x) = 0$ for all $\lambda \neq 0$. Since $k$ is infinite, then $G$ is infinitely many roots, and hence, $G = 0$. It follows that all the coefficients defining $G$ are zero, giving us $F_i(x) = 0$. Since $x$ was an arbitrary choice of coordinates for $P$, then $F_i(x_1, ..., x_{n+1}) = 0$ for all $i$ and for all choice of coordinates for $P$. \\
\end{solution}

\begin{exercise}
    (a) Show that the definitions of this section carry over without change to the case where $k$ is an arbitrary field. (b) If $k_0$ is a subfield of $k$, show that $\P^n(k_0)$ may be identified with a subset of $\P^n(k)$. \\
\end{exercise}

\begin{solution}
    \begin{enumerate}
        \item The only definitions in this section are the definitions of $\P^n(k)$, $U_i$ and $H_{\infty}$. None of these definitions used (explicitly or implicitly) properties not shared by all fields. Thus, they carry over without change to the case where $k$ is an arbitrary field.
        \item Consider the function $i : \P^n(k_0) \to \P^n(k)$ defined by
        $$i([x_1: ... : x_{n+1}]_{\P^n(k_0)}) = [x_1: ... : x_{n+1}]_{\P^n(k)}$$
        where the subscript denotes the space in which the point lives in. First, let's show that this function is well-defined by showing that the image of a point doesn't depend on the choice of coordinates. Let $P  = [x_1 : ... : x_{n+1}]_{\P^n(k_0)}$, then any other choice of coordinates for $P$ would be of the form $[\lambda x_1 : ... : \lambda x_{n+1}]_{\P^n(k_0)}$ for some non-zero $\lambda \in k_0$. Hence:
        \begin{align*}
            i([\lambda x_1 : ... : \lambda x_{n+1}]_{\P^n(k_0)}) &= [\lambda x_1 : ... : \lambda x_{n+1}]_{\P^n(k)} \\
            &= [x_1 : ... : x_{n+1}]_{\P^n(k)} \\
            &= i([x_1 : ... : x_{n+1}]_{\P^n(k_0)}).
        \end{align*}
        Therefore, $i$ is well-defined. Now, let's show that $i$ is injective. Let $P = [x_1 : ... : x_{n+1}]_{\P^n(k_0)}$ and $Q = [y_1 : ... : y_{n+1}]_{\P^n(k_0)}$ such that $i(P) = i(Q)$, then equivalently:
        $$[x_1 : ... : x_{n+1}]_{\P^n(k)} = [y_1 : ... : y_{n+1}]_{\P^n(k)}.$$
        This implies that $y_i = \lambda x_i$ for all $i$ and for some non-zero constant $\lambda \in k$. Let $i_0$ be such that $x_{i_0} \neq 0$, then it follows that $\lambda = y_{i_0}/x_{i_0} \in k_0$. Thus, $y_i = \lambda x_i$ for all $i$ and for some non-zero constant $\lambda \in k_0$. It follows that
        $$P = [x_1 : ... : x_{n+1}]_{\P^n(k_0)} = [y_1 : ... : y_{n+1}]_{\P^n(k_0)} = Q.$$
        Therefore, $i$ is injective, and hence, we can naturally identify $\P^n(k_0)$ with a subset of $\P^n(k)$ (the image of $i$).
    \end{enumerate}
\end{solution}
\section{Projective Algebraic Sets}

\begin{exercise}
    Let $I$ be a homogeneous ideal in $k[X_1, ..., X_{n+1}]$. Show that $I$ is prime if and only if the following condition is satisfied; for any forms $FG \in k[X_1, ..., X_{n+1}]$, if $FG \in I$, then $F\in I$ or $G \in I$. \\
\end{exercise}

\begin{solution}
    If $I$ is prime, then the condition is clearly satisfied. Suppose now that $I$ satisfies the given condition and let $F,G \in k[X_1, ..., X_{n+1}]$ such that $F = \sum F_i$, $G = \sum G_i$ and $FG \in I$. Since
    $$FG = F_0G_0 + (F_0G_1 + F_1G_0) + (F_0G_2 + F_1G_1 + F_2G_0) + ...,$$
    and $I$ is homogeneous, then each term in the sum above is in $I$. Hence, $F_0G_0 \in I$. Since $I$ satisfies the given condition, then $F_0 \in I$ or $G_0 \in I$. If $F_0 \notin I$, then $G_0 \in I$, and so
    $$F_0G_1 = (F_0G_1 + F_1G_0) - F_1G_0 \in I.$$
    This again implies that $G_1 \in I$ since $F_0 \in I$. By induction, we can show that this case implies that $G_i \in I$ for all $i$, and hence, that $G \in I$. Similarly, if $F_0 \in I$, then
    $$F_1G_0 = (F_0G_1 + F_1G_0) - F_0G_1 \in I$$
    and hence, $F_1 \in I$ or $G_0 \in I$. by repeating this procedure, we can either fall in the case where all the $F_i$'s are in $I$, and hence deduce that $F \in I$, or, there will be one $F_i$ which is not in $I$, and hence, forcing all the $G_i$'s to be in $I$, which would imply that $G \in I$. Thus, in all possible cases, $F \in I$ or $G \in I$. Therefore, $I$ is prime. \\
\end{solution}

\begin{exercise}
    If $I$ is a homogeneous ideal, show that $\rad(I)$ is also homogeneous. \\
\end{exercise}

\begin{solution}
    Let $F \in \rad(I)$ such that $F = \sum_{i=m}^d F_i$ where $\deg(F_i) = i$, then $F^n \in I$ for some $n \geq 1$. It follows that
    $$F^n = F_m^n + [\text{terms of higher degree}] \in I$$
    Since $I$ is homogeneous and $F_m^n$ is the form of smallest degree in the decomposition of $F^n$ into a sum of forms, then $F_m^n \in I$. Hence, $F_m \in \rad(I)$, and so $F - F_m \in \rad(I)$. It suffices to repeat this argument (induction) to get that $F_i \in \rad(I)$ for all $i$. Therefore, $\rad(I)$ is homogeneous. \\
\end{solution}

\begin{exercise}
    State and prove the projective analogues of properties (1)-(10) of Chapter 1, Sections 2 and 3. \\
\end{exercise}

\begin{solution}
    \begin{enumerate}[label = (\arabic*)]
        \item Let's show that if $I$ is the ideal in $k[X_1, ..., X_{n+1}]$ generated by the subset $S \subset k[X_1, ..., X_{n+1}]$, then $V(S) = V(I)$. Let $P \in V(I)$, then $F(P) = 0$ for all $F \in I$. Since $S \subset I$, then $F(P) = 0$ for all $F \in S$. Hence, $V(I) \subset V(S)$. Conversely, let $P \in V(S)$, then $F(P) = 0$ for all $F \in S$. Let $G \in I$, then $G = \sum \alpha_i F_i$ for some $\alpha_i \in k[X_1, ..., X_{n+1}]$ and $F_i \in S$. Let $(x_1, ..., x_{n+1})$ be a choice of homogeneous coordinates for $P$, then
        $$G(x_1, ..., x_{n+1}) = \sum \alpha_i(x_1, ..., x_{n+1}) F_i(x_1, ..., x_{n+1}) = 0.$$
        Since it holds for all choice of homogeneous coordinates for $P$, then $G(P) = 0$. Since it holds for all $G \in I$, then $P \in V(I)$. Therefore, $V(S) = V(I)$.
        \item Let's show that if $\{I_{\alpha}\}$ is any collection of ideals, then $V(\bigcup_{\alpha}I_{\alpha}) = \bigcap_{\alpha}V(I_{\alpha})$. Let $P \in V(\bigcup_{\alpha}I_{\alpha})$, then $F(P) = 0$ for all $F \in \bigcup_{\alpha}I_{\alpha}$. Let $\alpha$ be arbitrary and $F \in I_{\alpha}$, then $F \in \bigcup_{\alpha}I_{\alpha}$, and hence, $F(P) = 0$. Thus, $P \in V(I_{\alpha})$. Since $\alpha$ was set to be arbitrary, then $P \in \bigcap_{\alpha}V(I_{\alpha})$. Conversely, let $P \in \bigcap_{\alpha}V(I_{\alpha})$ and take any $F \in \bigcup_{\alpha}I_{\alpha}$, then $F \in I_{\alpha'}$ for some $\alpha'$. Since $P \in V(I_{\alpha'})$, then $F(P) = 0$. Since it holds for all $F \in \bigcup_{\alpha}I_{\alpha}$, then $P \in V(\bigcup_{\alpha}I_{\alpha})$. Therefore, $V(\bigcup_{\alpha}I_{\alpha}) = \bigcap_{\alpha}V(I_{\alpha})$.
        \item Let's show that if $I \subset J$, then $V(I) \supset V(J)$. Let $P \in V(J)$, then $F(P) = 0$ for all $F \in J$. Let $F \in I$, then $F \in J$, and hence, $F(P) = 0$. Since it holds for all $F \in I$, then $P \in V(I)$. Therefore, $V(I) \supset V(J)$.
        \item Let's show that $V(FG) = V(F)\cup V(G)$, and that $V(I) \cup V(J) = V(\{FG | F\in I, G \in J\})$. Let $P \in V(FG)$, then $(FG)(P) = 0$. Hence, $F(x_1, ..., x_{n+1}) = 0$ or $G(x_1, ..., x_{n+1}) = 0$ for all choice of homogeneous coordinates for $P$. Fix a choice of homogeneous coordinates $P = [x_1 : ... : x_{n+1}]$ and define the polynomials
        $$f(\lambda) = F(\lambda x_1, ..., \lambda x_{n+1}), \qquad g(\lambda) = G(\lambda x_1, ..., \lambda x_{n+1}).$$
        By our previous remark, for all non-zero $\lambda \in k$, either $f(\lambda) = 0$ or $g(\lambda) = 0$. Since $k$ is infinite, then it follows that either $f$ or $g$ has infinitely many zeros, without loss of generality, suppose that $f$ has infinitely many zeros. Since $f$ is a one-variable polynomial, then we must have $f = 0$. Equivalently, this means that $F(P) = 0$. Thus, $P \in V(F) \subset V(F) \cup V(G)$. Conversely, suppose that $P \in V(F)\cup V(G)$ and suppose that $P \in V(F)$ without loss of generality, then $F(P) = 0$. Equivalently, this means that $F(x_1, ..., x_{n+1}) = 0$ for all choice of homogeneous coordinates for $P$. Hence, $F(x_1, ..., x_{n+1})G(x_1, ..., x_{n+1}) = 0$ for all choice of coordinates for $P$. By definition, this implies that $(FG)(P) = 0$, and hence, that $P \in V(FG)$. Therefore, $V(FG) = V(F) \cup V(G)$.\\
        
        Let's now prove the more general statement using what we just proved. Let $P \in V(\{FG | F \in I, G \in J\})$, then $(FG)(P) = 0$ for all $F \in I$ and $G \in J$. By the previous proof, it follows that $F(P) = 0$, or $G(P) = 0$. Now, suppose that $P \in V(I)$, then we are done. Otherwise, there is a $F_0 \in I$ such that $F_0(P) \neq 0$. Thus, for all $G \in J$, the equation $(F_0G)(P) = 0$ implies that $F_0(P) = 0$ or $G(P) = 0$, which implies that $G(P) = 0$. Thus, in that case, $P \in V(J)$. Thus, in both cases, $P \in V(I) \cup V(J)$. Conversely, let $P \in V(I) \cup V(J)$ and suppose without loss of generality that $P \in V(I)$, then $F(P) = 0$ for all $F \in I$. It follows that $(FG)(P) = 0$ for all $F \in I$ and $G \in J$. Therefore, $P \in V(\{FG | F \in I, G \in J\})$. Therefore, $V(I) \cup V(J) = V(\{FG | F\in I, G \in J\})$.
        \item Let's show that $V(0) = \P^n$; $V(1) = \varnothing$; $\{[x_1 : ... : x_{n+1}]\}$ is a projective algebraic set for all $P \in \P^n$. First, we clearly have that every point in $\P^n$ is a zero of the zero polynomial so $V(0) = \P^n$. Similarly, no point in $\P^n$ is a zero of the polynomial 1 so $V(1) = \varnothing$. Next, let $P$ be a point in $\P^n$ and fix some homogeneous coordinates $(x_1, ..., x_{n+1})$ for $P$. In $\A^{n+1}(k)$, the line $L$ through the origin and $(x_1, ..., x_{n+1})$ is an algebraic set, so it can be described as $V(f_1, ..., f_r) \subset \A^{n+1}(k)$ for some polynomials $f_i \in k[x_1, ..., x_{n+1}]$ (Problem 2.15). For all $\lambda \neq 0$, the point $(\lambda x_1, ..., \lambda x_{n+1})$ is in the line, and hence, it is in $V(f_1, ..., f_r)$. Thus, $f_i(\lambda x_1, ..., \lambda x_{n+1}) = 0$ for all $i$ and for all $\lambda \neq 0$. By definition, this means that $P$ is a zero for all $f_i$, so $\{P\} \subset V(f_1, ..., f_r)$ viewed now as a subset of $\P^n$. Conversely, let $Q \in \P^n$ such that $Q \in V(f_1, ..., f_r)$, then choosing some coordinates $(y_1, ..., y_{n+1})$ for $Q$, we get that $(y_1, ..., y_{n+1}) \in V(f_1, ..., f_r) = L \subset \A^{n+1}(k)$. Hence, $(y_1, ..., y_{n+1}) = (\lambda x_1, ..., \lambda x_{n+1})$ for some $\lambda \neq 0$ so $Q = P$. Therefore, $V(f_1, ..., f_r) = \{P\}$ where $V(f_1, ..., f_r)$ is now viewed as a subset of $\P^n$.
        \item Let's prove that if $X \subset Y$, then $I(X) \supset I(Y)$. Let $F \in I(Y)$, then $F(P) = 0$ for all $P \in Y$. Let $P \in X$, then $P \in Y$, and hence, $F(P) = 0$. Since it holds for all $P \in X$, then $F \in I(X)$. Therefore, $I(X) \supset I(Y)$.
        \item Let's show that $I(\varnothing) = k[X_1, ..., X_{n+1}]$ and $I(\P^n(k)) = (0)$. First, for all $F \in k[X_1, ..., X_{n+1}]$, all the points in $\varnothing$ are zeros of $F$, so $F \in I(\varnothing)$. Therefore, $I(\varnothing) = k[X_1, ..., X_{n+1}]$. Next, if $F \in I(\P^n(k))$, then $F(\A^{n+1}(k)\setminus \{0\}) = \{0\}$. If we define $f(\lambda) = F(\lambda, \dots, \lambda)$, we get that $f$ is zero for all $\lambda \neq 0$. Since $f$ is a one-variable polynomial and $k$ is infinite, then $f = 0$, and so $F(0) = 0$. Therefore, $F \in I(\A^{n+1}(k))$, and hence, $F = 0$. It follows that $I(\P^n(k)) = (0)$.
        \item Let's show that $I(V(S)) \supset S$ for any set $S$ of polynomials; $V(I(X)) \supset X$ for any set $X$ of points. Let $F \in S$ and $P \in V(S)$, then by definition of $V(S)$, $F(P) = 0$. Since it holds for all $P \in V(S)$, then $F \in I(V(S))$. Therefore, $I(V(S)) \supset S$. Next, Let $P \in X$ and $F \in I(X)$, then by definition, $F(P) = 0$. Since it holds for all $F \in I(X)$, then $P \in V(I(X))$. Therefore, $V(I(X)) \supset X$.
        \item Let's show that
        $$V(I(V(S))) = V(S) \quad \text{ and } \quad I(V(I(X))) = I(X)$$
        for any set $S$ of polynomials, and any set $X$ of points. First, we know that $V(I(V(S))) \supset V(S)$ and $I(V(S)) \supset S$ from property (8). The inclusion $I(V(S)) \supset S$ implies $V(I(V(S))) \subset V(S)$ using property (3). Therefore, $V(I(V(S))) = V(S)$. Next, we have the inclusions $I(V(I(X))) \supset I(X)$ and $V(I(X)) \supset X$ from property (8). The latter implies that $I(V(I(X))) \subset I(X)$ using property (6). Therefore, $I(V(I(X))) = I(X)$.
        \item Let's show that $I(X)$ is a radical ideal for any $X \subset \P^n(k)$. Let $F^n \in I(X)$, then $F^n(P) = 0$ for all $P \in X$. By property (4), it follows that $F(P) = 0$ or $F^{n-1}(P) = 0$. If $F^{n-1}(P) = 0$, simply reapply the same procedure to get that in all cases, $F(P) = 0$. Since it holds for all $P \in X$, then $F \in I(X)$. Therefore, $I(X)$ is a radical ideal. \\
    \end{enumerate}
\end{solution}

\begin{exercise}
    Show that each irreducible component of a cone is also a cone. \\
\end{exercise}

\begin{solution}
    From the equation $C(V_p(I)) = V_a(I)$, we get that the cone of a projective algebraic subset of $\P^n(k)$ is an affine algebraic set of $\A^{n+1}(k)$. Morever, it is easy to show that $C(V_1 \cup V_2) = C(V_1) \cup C(V_1)$. Hence, if we take a projective algebraic set $V$ such that $V = V_1 \cup \dots \cup V_r$ is its unique decomposition into irreducible components, then $C(V) = C(V_1) \cup \dots \cup C(V_r)$. Let's show that $C(V_1), ..., C(V_r)$ are the irreducible components of $C(V)$. First, they are all irreducible because $I_a(C(V_i)) = I_p(V_i)$ is prime. Morever, suppose that $C(V_i) \subset C(V_j)$, then $I_p(V_i) \subset I_p(V_j)$, and hence, $V_i \subset V_j$. But this is a contradiction since $V_i$ and $V_j$ are irreducible components of $V$. Therefore, $C(V_1), ..., C(V_r)$ are the irreducible components of $C(V)$. By uniqueness of the decomposition into irreducible components, we get that each irreducible component of a cone is also a cone. \\
\end{solution}

\begin{exercise}
    Let $V = \P^1$, $\Gamma_h(V) = k[X,Y]$. Let $t = X/Y \in k_h(V)$, and show that $k(V) = k(t)$. Show that there is a natural one-to-one correspondence between the points of $\P^1$ and the DVR's with quotient field $k(V)$ that contain $k$ (see Problem 2.27); which DVR corresponds to the point at infinity? \\
\end{exercise}

\begin{solution}
    Let $f$ and $g$ be two forms in $\Gamma_h(V)$ of the same degree, say $d$, then we can write $f = \sum_{i=0}^{d}a_i X^{d-i}Y^i$ and $g = \sum_{i=0}^{d}b_iX^{n-i}Y^i$ where $a_i, b_i \in k$, then
    \begin{align*}
        \frac{f}{g} &= \frac{\sum_{i=0}^{d}a_i X^{d-i}Y^i}{\sum_{i=0}^{d}b_i X^{d-i}Y^i} \\
        &= \sum_{i=0}^{d}a_i\frac{X^{d-i}Y^i}{\sum_{j=0}^{d}b_j X^{d-j}Y^j} \\
        &= \sum_{i=0}^{d}a_i\left(\frac{\sum_{j=0}^{d}b_j X^{d-j}Y^j}{X^{d-i}Y^i}\right)^{-1} \\
        &= \sum_{i=0}^{d}a_i\left(\sum_{j=0}^{d}b_j\frac{X^{d-j}Y^j}{X^{d-i}Y^i}\right)^{-1} \\
        &= \sum_{i=0}^{d}a_i\left(\sum_{j=0}^{d}b_j\frac{X^{i-j}}{Y^{i-j}}\right)^{-1} \\
        &= \sum_{i=0}^{d}a_i\left(\sum_{j=0}^{d}b_j t^{i-j}\right)^{-1} \in k(t).\\
    \end{align*}
    Hence, $k(V) \subset k(t)$. Conversely, since $k$ and $t$ are in $k(V)$ and $k(V)$ is a field, then $k(t) \subset k(V)$. Therefore, $k(V) = k(t)$. 

    Next, since $k(V) = k(t)$, then the DVR's with quotient field $k(V)$ that contain $k$ are precisely $\O_a$ for $a \in \A^1(k)$ with uniformizing parameter $t = X - a$, and $\O_{\infty}$ with uniformizing parameter $t = 1/X$. Associate each point $[a : 1] \in \P$ to $\O_a$ and the point at infinity $[0 : 1]$ to $\O_{\infty}$. \\
\end{solution}

\begin{exercise}
    Let $I$ be a homogeneous ideal in $k[X_1, ..., X_{n+1}]$, and
    $$\Gamma = k[X_1, ..., X_{n+1}]/I.$$
    Show that the forms of degree $d$ in $\Gamma$ form a finite-dimensional vector space over $k$. \\
\end{exercise}

\begin{solution}
    \td \\
\end{solution}

\begin{exercise}
    \td \\
\end{exercise}

\begin{solution}
    \td \\
\end{solution}

\begin{exercise}
    \td \\
\end{exercise}

\begin{solution}
    \td \\
\end{solution}

\begin{exercise}
    Let $H_1, ..., H_m$ be hyperplanes in $\P^n$, $m \leq n$. Show that $H_1 \cap H_2 \cap ... \cap H_m \neq \varnothing$. \\
\end{exercise}

\begin{solution}
    Since $H_i$ is a hyperplane, then $H_i = V(f_i)$ where $f_i \in k[X_1, ..., X_{n+1}]$ is a form of degree. Hence, $\bigcap_i H_i = V(f_1, ..., f_m)$. An element $P \in \P^n$ is an element of $\bigcap_i H_i$ if and only if $f_i(P) = 0$ for all $1 \leq i \leq m$. If we write $f_i = a_{i1}X_1 + a_{i2}X_2 + \dots + a_{i(n+1)}X_{n+1}$, then the condition is equivalent to the solvability of the following system of equations:
    \begin{align*}
        a_{11}x_1 + a_{12}x_2 + \dots + a_{1(n+1)}x_{n+1} &= 0, \\
        a_{21}x_1 + a_{22}x_2 + \dots + a_{2(n+1)}x_{n+1} &= 0, \\
        \vdots \qquad \qquad \qquad \qquad \vdots \qquad & \quad \ \vdots \\
        a_{m1}x_1 + a_{m2}x_2 + \dots + a_{m(n+1)}x_{n+1} &= 0.
    \end{align*}
    Define the linear transformation $T : k^{n+1} \to k^m$ as follows:
    $$T(\vec{x}) = \begin{pmatrix}
        a_{11} & a_{12} & \dots & a_{1(n+1)} \\
        a_{21} & a_{22} & \dots & a_{2(n+1)} \\
        \ldots & \ldots & \ddots & \ldots \\
        a_{m1} & a_{m2} & \dots & a_{m(n+1)}
    \end{pmatrix}\vec{x},$$
    then the system of equation is equivalent to the equation $T\vec{x} = 0$ with $x \in k^{n+1}$. By the rank-nullity theorem, we have that
    $$\dim \ker T + \dim \Ima T = n+1.$$
    Since $\Ima T \subset k^m$, then
    $$\dim \ker T = (n+1) - \dim \Ima T \geq (n+1) - m \geq 1.$$
    Thus, there is at least one non-zero vector $(x_1, ..., x_{n+1})$ such that all of its scalar multiples are solution to the equation (by linearity). Hence, if we define $P = [x_1 : ... : x_{n+1}] \in \P^n$, then another way to phrase the previous sentence is to say that $f_i(P) = 0$ for all $i$, or equivalently, that $P \in \bigcap_i H_i$. Therefore, $H_1 \cap H_2 \cap ... \cap H_m \neq \varnothing$. \\
\end{solution}

\begin{exercise}
    \td \\
\end{exercise}

\begin{solution}
    \td \\
\end{solution}

\begin{exercise}
    \td \\
\end{exercise}

\begin{solution}
    \td \\
\end{solution}

\begin{exercise}
    \td \\
\end{exercise}

\begin{solution}
    \td \\
\end{solution}

\begin{exercise}
    \td \\
\end{exercise}

\begin{solution}
    \td \\
\end{solution}

\begin{exercise}
    \td \\
\end{exercise}

\begin{solution}
    \td \\
\end{solution}

\begin{exercise}
    \td \\
\end{exercise}

\begin{solution}
    \td \\
\end{solution}


\section{Affine and Projective Varieties}

\begin{exercise}
    If $I = (F)$ is the ideal of an affine hypersurface, show that $I^* = (F^*)$. \\
\end{exercise}

\begin{solution}
    Since $F \in I$, then $F^* \in I^*$ and hence, $(F^*) \subset I^*$. Conversely, let $h F \in (F)$, then $(hF)^*$ is one of the generators of $I^*$. But notice that $(hF)^* = h^*F^* \in (F^*)$ so all the generators of $I^*$ are in $(F^*)$. Therefore, $I^* \subset (F^*)$, and hence, $I^* = (F^*)$. \\
\end{solution}

\begin{exercise}
    Let $V = V(Y-X^2, Z - X^3) \subset \A^3$. Prove:
    \begin{enumerate}
        \item $I(V) = (Y - X^2, Z - X^3)$.
        \item $ZW - XY \in I(V)^*\subset k[X,Y,Z,W]$, but $ZW-XY \notin ((Y-X^2)^*, (Z-X^3)^*)$.
    \end{enumerate}
    So if $I(V) = (F_1, ..., F_r)$, it does not follow that $I(V)^* = (F_1^*, ..., F_r^*)$. \\
\end{exercise}

\begin{solution}
    \begin{enumerate}
        \item The ideal $(Y - X^2, Z - X^3)$ is prime because $k[X,Y,Z]/(Y - X^2, Z - X^3) \cong k[X]$. Hence, by the Nullstellenzats, $I(V) = (Y - X^2, Z - X^3)$.
        \item Since $Y - X^2, Z - X^3 \in I(V)$, then
        $$Z - XY = (Z - X^3) - X(Y - X^2) \in I(V).$$
        It follows that
        $$WZ - XY = (Z-XY)^* \in I(V)^*.$$
        Next, suppose $WZ - XY \in ((Y-X^2)^*, (Z-X^3)^*) = (WY-X^2, W^2Z-X^3)$, then there exist polynomials $A, B \in k[X,Y,Z,W]$ such that
        $$WZ - XY = A(X,Y,Z,W)(WY-X^2) + B(X,Y,Z,W)(W^2Z-X^3).$$
        Taking $X = Y = 0$ gives us 
        $$WZ = B(0,0,Z,W)W^2Z$$
        which is equivalent to
        $$1 = B(0,0,Z,W)W.$$
        But this last equation is impossible so $WZ - XY \notin ((Y-X^2)^*, (Z-X^3)^*)$. \\
    \end{enumerate}
\end{solution}

\begin{exercise}
    Show that if $V \subset W \subset \P^n$ are varieties, and $V$ is a hypersurface, then $W = V$ or $W = \P^n$ (see Problem 1.30). \\
\end{exercise}

\begin{solution}
    Since $V$ is a hypersurface and a variety, then there is an irreducible polynomial $f \in k[X_1, ..., X_{n+1}]$ such that $V(f)$. From the inclusion $V \subset W$, we get $I(W) \subset I(V(f)) = (f)$ where the last equality comes from the Nullstellenzats. Since $W$ is a variety, then $I(W)$ must also be a prime ideal. Suppose that $I(W)$ is not the zero ideal, then any non-zero element $h_0 \in I(W)$ is of the form $h_0 = h_1f$ where $h_1 \in k[X_1, ..., X_{n+1}]$. Since $I(W)$ is prime, then $h_1 \in I(W)$ or $f \in I(W)$. If $f \notin I(W)$, then we can construct the infinite sequence $h_n \in I(W)$ such that $h_0 = h_n f^n$. But this is impossible because it would mean that the degree of $h$ is infinite. Therefore, $f \in I(W)$, and hence, $I(V) = (f) \subset I(W)$. In that case, $I(W) = I(V)$. In other words, either $I(W) = (0)$, or $I(W) = I(V)$. In the first case, $W = \P^n$, and in the second case, $W = V$. \\
\end{solution}

\begin{exercise}
    Suppose $V$ is a variety in $\P^n$ and $V \supset H_{\infty}$. Show that $V = \P^n$ or $V = H_{\infty}$. If $V = \P^n$, $V_* = \A^n$, while if $V = H_{\infty}$, $V_* = \varnothing$. \\
\end{exercise}

\begin{solution}
    First, notice that we can write $H_{\infty}$ as $V(X_{n+1})$, which now makes it clear that $H_{\infty}$ is an irreducible hypersurface. Hence, from the inclusion $H_{\infty} \subset V \subset \P^n$, we get $V = \P^N$ or $V = H_{\infty}$ using Problem 4.21. If $V = \P^n$, then $I(V) = (0)$, so $I(V)_* = (0)$, and hence, $V_* = V(I(V)_*) = V(0) = \A^n$. Similarly, if $V = H_{\infty}$, then $I(V) = (X_{n+1})$, so $I(V)_* = (1)$, and hence, $V_* = V(I(V)_*) = V(1) = \varnothing$. \\
\end{solution}

\begin{exercise}
    Describe all subvarieties in $\P^1$ and in $\P^2$. \\
\end{exercise}

\begin{solution}
    First, we know that all subvarieties of $\A^n$ can be lifted to subvarieties of $\P^n$ through the correspondence $V \mapsto V^*$. The subvarieties of $\A^1$ are the empty set, the singeltons and $\A^1$. It follows that the empty set, the singeltons of $U_2$ and $\P^n$ are subvarieties of $\P^n$. Since all singeltons are subvarieties, then we can add to this list the point at infinity. Finally, any subvariety that is contained in $H_{\infty}$ is either $\varnothing$ or $H_{\infty}$ (since $H_\infty$ is a singelton in this case) and any subvariety that contains $H_{\infty}$ is either equal to $H_{\infty}$ or $\P^n$ by the previous problem. Therefore, the complete list of subvarieties of $\P^1$ is: $\varnothing$, $\{P\}$ with $P \in \P^n$, and $\P^n$. \\
    
    Let's find the subvarieties of $\P^2$ in a similar way. First, recall that the subvarieties of $\A^2$ are: $\varnothing$, points, $V(f)$ where $f$ is irreducible and non-constant, and $\A^2$. It follows that $\varnothing$, points (including the points in $H_{\infty}$ since we already know that all singeltons are algebraic sets), $V(f^*)$ where $f$ is an irreducible non-constant polynomial, and $\P^n$ are subvarieties of $\P^2$. The last potential subvarieties of $\P^2$ are the ones that are either contained in $H_{\infty}$, or that contains $H_{\infty}$. In the first case, by the correspondence between $H_{\infty}$ and $\P^1$, we get that these varieties must be the empty set, the points in $H_{\infty}$ or $H_{\infty}$ entirely. In the second case, the subvarieties must be either $H_{\infty}$ or $\P^n$ (using Problem 4.22 again). Thus, in our list of subvarieties of $\P^n$, we only need to add $H_{\infty}$ to complete it. Therefore, the complete list of subvarieties of $\P^2$ is: $\varnothing$, all points, $V(f^*)$ where $f$ is an irreducible non-constant polynomial, $H_{\infty}$, and $\P^n$. \\
\end{solution}

\begin{exercise}
    \td \\
\end{exercise}

\begin{solution}
    \td \\
\end{solution}

\begin{exercise}
    \td \\
\end{exercise}

\begin{solution}
    \td \\
\end{solution}


\section{Multiprojective Space}

\begin{exercise}
    \td \\
\end{exercise}

\begin{solution}
    \td \\
\end{solution}

\begin{exercise}
    \td \\
\end{exercise}

\begin{solution}
    \td \\
\end{solution}

\begin{exercise}
    \td \\
\end{exercise}

\begin{solution}
    \td \\
\end{solution}



%CHAPTER 5
\chapter{Projective Plane Curves}

\td

%CHAPTER 6
\chapter{Varieties, Morphisms, and Rational Maps}

\section{The Zariski Topology}

\begin{exercise}
    Let $Z \subset Y \subset X$, $X$ a topological space. Give $Y$ the induced topology. Show that the topology induced by $Y$ on $Z$ is the same as that induced by $X$ on $Z$. \\
\end{exercise}

\begin{solution}
    Let $U \subset Z$ be an open set from the topology induced by $Y$, then there is an open subset $V$ of $Y$ such that $U = Z \cap V$. Similarly, since $V$ is from the topology induced by $X$ on $Y$, there is an open subset $W$ of $X$ such that $V = Y \cap W$. Hence, $U = Z \cap Y \cap W = Z \cap W$, from which it follows that $U$ is also open in the sense of the topology induced by $X$ on $Z$. 
    
    Conversely, let $U \subset Z$ be an open set from the topology induced by $X$ on $Z$, then there is an open subset $V$ of $X$ such that $U = Z \cap V$. Thus, we can write $U = Z \cap (Y \cap V)$. Since $Y \cap V$ is an open subset of $Y$ from the topology induced by $X$ on $Y$, then $U = Z \cap (Y \cap V)$ is an open subset from the topology induced by $Y$ on $Z$. Therefore both topologies are the same. \\
\end{solution}

\begin{exercise}
    (a) Let $X$ be a topological space, $X = \bigcup_{\alpha \in \mathcal{A}} U_{\alpha}$, $U_{\alpha}$ open in $X$. Show that a subset $W$ of $X$ is closed if and only if each $W \cap U_{\alpha}$ is closed (in the induced topology) in $U_{\alpha}$. (b) Suppose similarly $Y = \bigcup_{\alpha\in \mathcal{A}}V_{\alpha}$, $V_{\alpha}$ open in $Y$, and suppose $f : X \to Y$ is mapping such that $f(U_{\alpha}) \subset V_{\alpha}$. Show that $f$ is continuous if and only if the restriction of $f$ to each $U_{\alpha}$ is a continuous mapping from $U_{\alpha}$ to $V_{\alpha}$. \\
\end{exercise}

\begin{solution}
    \begin{enumerate}
        \item Suppose $W$ is closed, then from the equation
        $$U_{\alpha}\setminus (W \cap U_{\alpha}) = U_{\alpha} \cap (X \setminus W),$$
        we get that $W \cap U_{\alpha}$ must be closed in $U_{\alpha}$ for all $\alpha$. Conversely, suppose that $W \cap U_{\alpha}$ is closed in $U_{\alpha}$ for all $\alpha$, then by the previous equation $U_{\alpha}\cap (X \setminus W)$ is an open subset of $X$ for all $\alpha$. It follows that
        $$\bigcup_{\alpha}\left[U_{\alpha}\cap (X \setminus W)\right] = (X \setminus W) \cap \bigcup_{\alpha}U_{\alpha} = (X \setminus W)\cap X = X \setminus W$$
        is open in $X$. Therefore, $W$ is a closed subset of $X$.
        \item Suppose that $f$ is continuous and let $U$ be an open subset of $V_{\alpha}$, then $U = V_{\alpha}\cap V$ where $V$ is an open subset of $Y$. It follows that
        $$f|_{U_{\alpha}}^{-1}(U) = U_{\alpha}\cap f^{-1}(V_{\alpha}\cap V).$$
        Since both $V_{\alpha}$ and $V$ are open subsets of $Y$, then $V_{\alpha}\cap V$ is also an open subset of $Y$. Thus by continuity of $f$, the set $f|_{U_{\alpha}}^{-1}(U)$ is open in $U_{\alpha}$ since $f^{-1}(V_{\alpha}\cap V)$ is open in $X$. Therefore, $f|_{U_{\alpha}} : U_{\alpha} \to V_{\alpha}$ is continuous for all $\alpha$. Conversely, suppose that $f_{U_{\alpha}}$ is continuous for all $\alpha$ and let $U$ be an open subset of $Y$, then
        $$f^{-1}(U) = \bigcup_{\alpha}U_{\alpha}\cap f^{-1}(U) = \bigcup_{\alpha}f|_{U_{\alpha}}^{-1}(U \cap V_{\alpha}).$$
        By our assumption, $f|_{U_{\alpha}}^{-1}(U \cap V_{\alpha})$ must be open since $U \cap V_{\alpha}$ is open in $V_{\alpha}$ and $f|_{U_{\alpha}}$ is continuous. Thus, $f^{-1}(U)$ is open. Therefore, $f$ is continuous on $X$. \\
    \end{enumerate}
\end{solution}

\begin{exercise}
    (a) Let $V$ be an affine variety, $f \in \Gamma(V)$. Considering $f$ as a mapping from $V$ to $k = \A^1$, show that $f$ is continuous. (b) Show that any polynomial map of affine varieties is continuous. \\
\end{exercise}

\begin{solution}
    \begin{enumerate}
        \item We can suppose that $f$ is the restriction of a polynomial to $V$. Recall that the algebraic subsets of $\A^1$ are the empty set, the finite sets, and $\A^1$. Let's show that in each case, $f^{-1}(C)$ is a closed subset of $V$ where $C$ is a closed subset of $\A^1$. Trivially, $f^{-1}(\varnothing) = \varnothing$ and $f^{-1}(\A^1) = V$, which are both closed subsets of $V$. Finally, let $C = \{a_1, ..., a_r\}$ be a finite subset of $\A^1$, then
        $$f^{-1}(C) = V\cap (V(f - a_1) \cup V(f- a_2) \cup ... \cup V(f-a_r)).$$
        Since each $V(f-a_i)$ is closed in $\A^n$, then $V(f - a_1) \cup ... \cup V(f-a_r)$ is a closed subset of $\A^n$. It follows that $f^{-1}(C)$ is a closed subset of $V$. Therefore, $f$ is a continuous function from $V$ to $k = \A^1$.
        \item Let $\varphi : V \to W$ be a polynomial map and let $C = W \cap V(f_1, ..., f_r)$ be an arbitrary closed subset of $W$, then
        $$\varphi^{-1}(C) = V \cap V(f_1 \circ \varphi, ..., f_r \circ \varphi)$$
        which is a closed subset of $V$. therefore, polynomial maps are continuous.\\
    \end{enumerate}
\end{solution}

\begin{exercise}
    Let $U_i \subset \P^n$, $\varphi_i : \A^n \to U_i$ as in Chapter 4. Give $U_i$ the topology induced from $\P^n$. (a) Show that $\varphi_i$ is a homeomorphism. (b) Show that a set $W \subset \P^n$ if and only if each $\varphi_i^{-1}(W)$ is closed in $\A^n$, $i=1,...,n+1$. (c) Show that if $V \subset \A^n$ is an affine variety, then the projective closure $V^*$ of $V$ is the closure of $\varphi_{n+1}(V)$ in $\P^n$. \\
\end{exercise}

\begin{solution}
    \begin{enumerate}
        \item \td
        \item \td
        \item \td \\
    \end{enumerate}
\end{solution}

\begin{exercise}
    \td \\
\end{exercise}

\begin{solution}
    \td \\
\end{solution}

\begin{exercise}
    \td \\
\end{exercise}

\begin{solution}
    \td \\
\end{solution}

\begin{exercise}
    Let $V$ be an affine variety, $f \in \Gamma(V)$. (a) Show that $V(f) = \{P \in V | f(P) = 0\}$ is a closed subset of $V$, and $V(f) \neq V$ unless $f = 0$. (b) Suppose $U$ is a dense subset of $V$ and $f(P) = 0$ for all $P \in U$. Then $f = 0$. \\
\end{exercise}

\begin{solution}
    \begin{enumerate}
        \item Let $F$ be a polynomial such that $\overline{F} = f$, then $V(f) = V \cap V(F)$, which is closed in $V$. Next, suppose $V(f) = V$, then $F(P) = 0$ for all $P \in V$. Hence, $F \in I(V)$ which implies that $f = \overline{F} = 0$ in $\Gamma(V)$.
        \item First, we can restate the statement of the question by writing that $U \subset V(f)$. Moreover, $V(f)$ is a closed set so the closure of $U$ is also a subset of $U$. But $U$ is dense so its closure is $V$. Hence, $V \subset V(f)$, so $V(f) = V$. Therefore, $f = 0$. \\
    \end{enumerate}
\end{solution}

\begin{exercise}
    \td \\
\end{exercise}

\begin{solution}
    \td \\
\end{solution}


\section{Varieties}

\begin{exercise}
    Let $X = \A^2\setminus \{(0,0)\}$, an open subvariety of $\A^2$. Show that $\Gamma(X) = \Gamma(\A^2) = k[X,Y]$. \\
\end{exercise}

\begin{solution}
    Clearly, we have that $k[X,Y] \subset \Gamma(X)$ since polynomials are defined everywhere on $X$. Conversely, let $z \in \Gamma(X)$, then $z \in k(X,Y)$ and it is defined everywhere except maybe at zero. Since $k[X,Y]$ is a UFD, then we can write $z = a/b$ in a unique way (up to units) with $a,b \in k[X,Y]$. Hence, we get that $b(P) \neq 0$ for all $P \neq 0$. Let's prove that this implies that $b(0,0) \neq 0$ as well. \td [Prouver que $b \equiv c$ sur $k^2 \setminus \{(0,0)\}$, donc $b - c$ est zero sur un sous ensemble ouvert de $k^2$, donc dense. Donc, $b$ est zero sur $k^2$.]\td \\
\end{solution}

\begin{exercise}
    Let $U$ be an open subvariety of a variety $X$, $Y$ a closed subvariety of $U$. Let $Z$ be the closure of $Y$ in $X$. Show that
    \begin{enumerate}
        \item $Z$ is a closed subvariety of $X$.
        \item $Y$ is an open subvariety of $Z$.\\
    \end{enumerate}
\end{exercise}

\begin{solution}
    \begin{enumerate}
        \item \td
        \item \td \\
    \end{enumerate}
\end{solution}

\begin{exercise}
    \td \\
\end{exercise}

\begin{solution}
    \td \\
\end{solution}

\begin{exercise}
    \td \\
\end{exercise}

\begin{solution}
    \td \\
\end{solution}



\end{document}